% !TeX root = ../../Book5.tex
% 纲要标题：## 第4章　双陪集、Mackey 理论与 Clifford 理论
% 纲要位置：工作记录/计划纲要.md，第 4136 行。


\chapter{双陪集、Mackey 理论与 Clifford 理论}
\label{b5:ch:04}
\BookChapterNumber{b5:ch:04}

\begin{chapterabstract}
本章研究诱导与限制怎样组合，以及不可约表示在正规子群上如何分解。双陪集把陪集作用分成轨道，Mackey 公式把每个轨道写成一个新的诱导表示，由此计算交织空间并判定诱导表示的不可约性。Clifford 理论进一步说明，正规子群上的不可约成分组成一个共轭轨道，整个表示可从惯性群上的等型分量恢复。最后用障碍余圈描述表示延拓的障碍，并在重数空间上建立射影表示。
\end{chapterabstract}

\begin{learninggoals}
\begin{itemize}
\item 能由陪集作用求双陪集、交集子群及每个分块的维数。
\item 写出 Mackey 分解的显式映射，说明两种伴随在交织空间公式中的作用。
\item 能计算正规子群上不可约成分的轨道、惯性群与重数，并应用 Clifford 对应。
\item 能构造障碍余圈，检验余圈恒等式，判断延拓障碍及其对重数空间的影响。
\end{itemize}
预备知识为有限群的陪集与正规子群、复表示的完全可约性和 Schur 引理，以及\chapref{b5:ch:03}中的诱导、限制、余诱导和 Frobenius 互反性。前两节的模分解允许任意基域；不可约性判据和后两节的结构定理在复数域上讨论。
\end{learninggoals}

取 \(G=S_3\)，令 \(H\) 为由一个换位生成的二阶子群。从 \(H\) 的平凡表示诱导到 \(G\)，得到以三个左陪集为基的置换表示；再限制回 \(H\) 时，其中一个陪集固定，另两个互换。原来一次诱导所得的空间，限制后便按陪集轨道分成两块。一般子群 \(H,K\le G\) 之间，轨道会由双陪集编号，每块又来自较小子群上的诱导，这就是 Mackey 分解\footnote{马凯（George Whitelaw Mackey，1916-2006），美国数学家，研究群表示与调和分析，建立诱导表示的重要方法。}的出发点。若 \(H\) 正规，还可追问一个不可约 \(G\)-表示限制到 \(H\) 后，不同不可约成分为何随共轭一起移动。

% !TeX root = ../../Book5.tex
% 纲要标题：### 4.1 双陪集
% 纲要位置：工作记录/计划纲要.md，第 3819 行。
% 纲要细目（写作范围）：
% * 双陪集分解
% * 诱导后限制的几何分块
% * 小群的实际计算


\section{双陪集}
\label{b5:sec:0401}

把一个 \(H\)-表示诱导到 \(G\)，再只观察子群 \(K\) 的作用时，原来由 \(G/H\) 标记的各个向量空间并不会任意混合。哪些空间能够互相到达，取决于 \(K\) 在陪集集上的轨道。双陪集正是这些轨道的另一种写法。

\subsection{双陪集分解}
\label{b5:subsec:0401:01}

\begin{definition}\label{b5:def:double-coset}
	设 \(G\) 为有限群，\(H,K\leq G\)，\(x\in G\)。集合
	\[
		KxH=\{kxh:k\in K,\ h\in H\}
	\]
	称为由 \(x\) 确定的 \term[shuangpeiji]{双陪集}[double coset]。所有这些集合组成的集合记为 \(K\backslash G/H\)。这个记号表示一个商集，一般没有商群结构。
\end{definition}
\symidx{\(K\backslash G/H\)：双陪集集}

关系 \(y=kxh\) 是等价关系：单位元给出自反性，等式 \(x=k^{-1}yh^{-1}\) 给出对称性，两次这样的表达式相接给出传递性。因此，不同双陪集互不相交，它们的并为 \(G\)。

\begin{proposition}\label{b5:prop:double-coset-orbits}
	设 \(G\) 为有限群，\(H,K\leq G\)，\(x\in G\)。在左陪集集 \(G/H\) 上令 \(K\) 左乘作用。包含 \(xH\) 的轨道为 \(KxH/H\)，其稳定子群为
	\[
		L_x=K\cap xHx^{-1}.
	\]
	映射 \(kL_x\mapsto kxH\) 给出 \(K\)-集合的同构
	\[
		K/L_x\longrightarrow KxH/H.
	\]
	特别地，
	\begin{equation}\label{b5:eq:double-coset-size}
		|KxH|=\frac{|K|\cdot|H|}{|K\cap xHx^{-1}|}.
	\end{equation}
\end{proposition}
\begin{proof}
	元素 \(k\in K\) 固定 \(xH\)，当且仅当 \(x^{-1}kx\in H\)，即 \(k\in L_x\)。而 \(k_1xH=k_2xH\) 当且仅当 \(k_2^{-1}k_1\in L_x\)，所以所给映射良定义、单射，并由轨道的定义而满射。它与左乘作用相容。轨道含有 \([K:L_x]\) 个左 \(H\)-陪集，每个陪集含有 \(|H|\) 个元素，得到公式。
\end{proof}

于是，选取双陪集代表元，就是在 \(K\) 作用于 \(G/H\) 的每个轨道中选一个点。即使 \(H=K\)，交集也仍须写成 \(H\cap xHx^{-1}\)，不能把共轭后的子群忽略。

\subsection{诱导后限制的几何分块}
\label{b5:subsec:0401:02}

设 \(k\) 为任意域，\(W\) 为有限维左 \(k[H]\)-模。诱导表示采用
\[
	\operatorname{Ind}_H^G W=k[G]\otimes_{k[H]}W.
\]
选取一组左陪集代表元 \(T\subseteq G\) 后，有向量空间直和分解
\[
	\operatorname{Ind}_H^G W=\bigoplus_{t\in T}t\otimes W.
\]
若 \(a\in K\)，则 \(a\) 把 \(t\otimes W\) 送到 \(atH\) 所对应的空间。因此，每一个 \(K\)-轨道内的全部空间之和都是 \(K\)-不变子空间。

\begin{proposition}\label{b5:prop:double-coset-blocks}
	设 \(G\) 为有限群，\(H,K\leq G\)，\(k\) 为域，\(W\) 为有限维左 \(k[H]\)-模。对每个双陪集 \(D=KxH\)，令 \(k[D]\) 为 \(k[G]\) 中以 \(D\) 为基的子空间。它是左 \(k[K]\)、右 \(k[H]\) 的双模，且有左 \(k[K]\)-模直和分解
	\[
		\operatorname{Res}_K^G\operatorname{Ind}_H^G W =\bigoplus_{D\in K\backslash G/H}k[D]\otimes_{k[H]}W.
	\]
	其中 \(D=KxH\) 所对应的分量的维数为
	\[
		[K:K\cap xHx^{-1}]\dim_k W.
	\]
\end{proposition}
\begin{proof}
	群元素基按互不相交的双陪集分组，得到 \(k[G]=\bigoplus_D k[D]\)。左右作用都保持各个 \(D\)，故这是双模的直和。对右 \(k[H]\)-模作张量积，直和的包含与投影仍互为分裂，得到所述直和。

	每个 \(D\) 是 \([K:K\cap xHx^{-1}]\) 个左 \(H\)-陪集之并。选取这些陪集的代表元后，\(k[D]\) 是同样秩的自由右 \(k[H]\)-模；张量积因而是同样多个 \(W\) 的直和。
\end{proof}

这里的双模结构只要求左乘 \(K\)、右乘 \(H\) 保持 \(D\)，不要求 \(k[D]\) 是子代数。例如在 \(S_3\) 中，取 \(H=K=\langle(12)\rangle\)、\(D=H(13)H\)，则 \((13)\in D\)，但 \((13)^2=1\notin D\)，所以 \(k[D]\) 对乘法不封闭。上述各块也只保证在限制到 \(K\) 后不变；取 \(W=k\) 为平凡表示时，\((13)\) 把 \(1\otimes1\) 从 \(H\) 所对应的块送到 \(H(13)H\) 所对应的块。

\ProseParagraphBreak

这个分块不要求特征不整除群阶。它来自群元素基的划分，与表示是否完全可约无关；下节将在每一块内重新识别一个诱导表示。

\subsection{小群的双陪集计算}
\label{b5:subsec:0401:03}

先取 \(G=S_3\)、\(H=K=\bigl\langle(12)\bigr\rangle\)。左陪集 \(gH\) 可由点 \(g(3)\) 标记，\(H\) 在三点集上的轨道为 \(\{3\}\) 与 \(\{1,2\}\)。因此双陪集代表元可取 \(1\) 和 \(x=(13)\)。对应的交集分别为 \(H\) 和
\[
	H\cap xHx^{-1}=\bigl\langle(12)\bigr\rangle\cap\bigl\langle(23)\bigr\rangle=\{1\}.
\]
两个双陪集的大小为 \(2\) 与 \(4\)，恰好覆盖 \(S_3\)。若 \(W=k\) 是平凡表示，则三维置换表示限制到 \(H\) 后分成一个固定点的平凡表示，以及两点自由轨道的正则表示。

\ProseParagraphBreak

若改在复数域上取 \(H\) 的符号表示 \(W=\mathbb Cw\)，即 \((12)w=-w\)，陪集轨道仍是上述两个，但每块内部的作用需要重新计算。取
\[
	e_1=1\otimes w,\qquad e_2=(13)\otimes w,\qquad e_3=(23)\otimes w.
\]
由 \((12)(13)=(23)(12)\) 和 \((12)(23)=(13)(12)\)，得到
\[
	(12)e_1=-e_1,\qquad (12)e_2=-e_3,\qquad (12)e_3=-e_2.
\]
第一块 \(\mathbb Ce_1\) 是符号表示；第二块在基 \(e_2,-e_3\) 下由 \((12)\) 交换两个基向量，所以仍是 \(H\) 的正则表示。在第二块中，\(e_2-e_3\) 生成平凡直线，\(e_2+e_3\) 生成符号直线，故
\[
	\operatorname{Res}_H^{S_3}\operatorname{Ind}_H^{S_3}W \cong \mathbf1_H\oplus W^{\oplus2}.
\]
同一双陪集划分给出了同样的块数与维数，却并不决定各块中的表示结构。

\ProseParagraphBreak

再取 \(G=S_4\)，让 \(H=K=S_{\{1,2,3\}}\) 固定 \(4\)。\(G/H\) 同样可识别为四点集，轨道为 \(\{4\}\) 与 \(\{1,2,3\}\)。代表元取 \(1\) 和 \(x=(34)\)，后者的交集为固定 \(3,4\) 的二阶子群 \(\bigl\langle(12)\bigr\rangle\)。所以双陪集大小为 \(6\) 与 \(18\)，诱导平凡表示的两块维数为 \(1\) 与 \(3\)。

\ProseParagraphBreak

这些计算提供了一种实际方法：先寻找容易理解的陪集作用，再算轨道和稳定子群。直接枚举 \(kxh\) 往往会重复产生同一个元素，而轨道方法同时给出分块数、代表元和每块大小。

% !TeX root = ../../Book5.tex
% 纲要标题：### 4.2 Mackey 分解
% 纲要位置：工作记录/计划纲要.md，第 4144 行。
% 纲要细目（写作范围）：
% * 双陪集公式
% * 交织算子空间的计算
% * 不可约性判据


\section{Mackey 分解}
\label{b5:sec:0402}

双陪集分块中的 \(x\otimes W\) 受交集 \(K\cap xHx^{-1}\) 作用。这个作用不是原封不动地把 \(W\) 限制到交集：必须先通过 \(x\) 把交集搬回 \(H\)。将这一点写清以后，Mackey 公式便可由一个显式张量映射得到。

\subsection{双陪集公式}
\label{b5:subsec:0402:01}

\begin{definition}\label{b5:def:conjugate-representation}
	设 \(G\) 为有限群，\(H\leq G\)，\(k\) 为域，\(\rho:H\to\operatorname{GL}_k(W)\) 为有限维表示，\(x\in G\)。在同一向量空间 \(W\) 上定义 \(xHx^{-1}\) 的作用
	\[
		\rho^x(a)=\rho(x^{-1}ax)\qquad(a\in xHx^{-1}).
	\]
	所得表示记为 \({}^xW\)，称为 \(W\) 的\term[gonge-biaoshi]{共轭表示}[conjugate representation]。
\end{definition}
\symidx{\({}^xW\)：共轭表示}

这个方向约定与左陪集相配合：在诱导表示中， \(a(x\otimes w)=x\otimes\rho(x^{-1}ax)w\)。此外，先取 \({}^yW\) 再取其 \(x\)-共轭，所得表示为 \({}^{xy}W\)。

\begin{theorem}[Mackey 分解]\label{b5:thm:mackey-decomposition}
	设 \(G\) 为有限群，\(H,K\leq G\)，\(k\) 为任意域，\(W\) 为有限维左 \(k[H]\)-模。取 \(K\backslash G/H\) 的一组代表元 \(X\)，并令 \(L_x=K\cap xHx^{-1}\)。则有左 \(k[K]\)-模同构
	\begin{equation}\label{b5:eq:mackey-decomposition}
		\operatorname{Res}_K^G\operatorname{Ind}_H^G W \cong \bigoplus_{x\in X} \operatorname{Ind}_{L_x}^K \operatorname{Res}_{L_x}^{xHx^{-1}}({}^xW).
	\end{equation}
\end{theorem}
这一分解称为 \term[mackey-fenjie]{Mackey 分解}[Mackey decomposition]，公式也称为 Mackey 双陪集公式。
\begin{proof}
	根据\cref{b5:prop:double-coset-blocks}，只须识别每一块。记 \(W\) 的 \(H\)-作用为 \(\rho\)。固定 \(x\)，把 \({}^xW\) 限制到 \(L_x\)，定义
	\[
		\begin{aligned}
			\Phi_x:k[K]\otimes_{k[L_x]}{}^xW &\longrightarrow k[KxH]\otimes_{k[H]}W,\\
			a\otimes w&\longmapsto ax\otimes w.
		\end{aligned}
	\]
	对 \(\ell\in L_x\)，有
	\[
		a\ell x\otimes w =ax(x^{-1}\ell x)\otimes w =ax\otimes\rho(x^{-1}\ell x)w.
	\]
	这正是源端平衡关系的像，故映射良定义；左 \(K\)-线性直接来自左乘。它的逆映射可以在群元素基上写成
	\[
		\Psi_x(k_0xh\otimes w)=k_0\otimes hw \qquad(k_0\in K,\ h\in H).
	\]
	这里须检查同一群元素的不同写法。如果 \(k_0xh=k_1xh_1\)，则
	\[
		\ell=k_1^{-1}k_0\in K\cap xHx^{-1}=L_x, \qquad (x^{-1}\ell x)h=h_1.
	\]
	源端的平衡关系因而给出
	\[
		k_0\otimes hw =k_1\ell\otimes hw =k_1\otimes\rho(x^{-1}\ell x)hw =k_1\otimes h_1w.
	\]
	所以 \(\Psi_x\) 与所选写法无关。对于 \(h_0\in H\)，它还把 \(k_0xhh_0\otimes w\) 与 \(k_0xh\otimes h_0w\) 都送到 \(k_0\otimes hh_0w\)，故经线性延拓后确实定义了张量积上的映射。直接代入可得 \(\Psi_x\Phi_x=1\)、\(\Phi_x\Psi_x=1\)。交集 \(L_x\) 正是同一元素写成 \(kxh\) 时产生的歧义所在，源端对 \(L_x\) 取张量积恰好消除了这种歧义。对所有 \(x\) 取直和即得结论。
\end{proof}

选取别的代表元会改变上式各项的具体写法，却不改变\cref{b5:prop:double-coset-blocks}中由双陪集本身确定的子空间。具体地，若 \(x'=k_0xh_0\)，其中 \(k_0\in K\)、\(h_0\in H\)，则 \(L_{x'}=k_0L_xk_0^{-1}\)，且对 \(\ell\in L_x\) 有
\[
	\begin{aligned}
		\rho^{x'}(k_0\ell k_0^{-1}) &=\rho(h_0^{-1}x^{-1}\ell xh_0),\\
		\rho(h_0)\rho^{x'}(k_0\ell k_0^{-1}) &=\rho^x(\ell)\rho(h_0).
	\end{aligned}
\]
后一等式使映射
\[
	k[K]\otimes_{k[L_{x'}]}{}^{x'}W \longrightarrow k[K]\otimes_{k[L_x]}{}^xW, \qquad a\otimes w\longmapsto ak_0\otimes\rho(h_0)w
\]
与平衡关系相容；其逆为 \(a\otimes w\mapsto ak_0^{-1}\otimes\rho(h_0^{-1})w\)。再经 \(\Phi_x\)，像为 \(ak_0x\otimes h_0w=ax'\otimes w\)，与 \(\Phi_{x'}\) 相同。因此两个代表元确实识别了同一个双陪集块。

\ProseParagraphBreak

回到 \(S_3\)、\(H=K=\langle(12)\rangle\) 的符号表示 \(W\)。代表元 \(1\) 的交集为 \(H\)，对应的项就是 \(W\)；代表元 \((13)\) 的交集为平凡群，对应的项是 \(\operatorname{Ind}_{\{1\}}^H\mathbb C=\mathbb C[H]\)。于是 Mackey 公式给出 \(W\oplus\mathbb C[H]\cong\mathbf1_H\oplus W^{\oplus2}\)，与上节直接计算的两块作用一致。

\subsection{交织算子空间的计算}
\label{b5:subsec:0402:02}

\begin{theorem}[交织空间公式]\label{b5:thm:mackey-intertwining}
	设 \(G\) 为有限群，\(H,K\leq G\)，\(k\) 为域，\(U,W\) 分别为有限维左 \(k[K]\)-模、左 \(k[H]\)-模。取双陪集集 \(K\backslash G/H\) 的代表元集 \(X\)，对 \(x\in X\) 令 \(L_x=K\cap xHx^{-1}\)。有 \(k\) 向量空间同构
	\begin{equation}\label{b5:eq:mackey-intertwining}
		\begin{aligned}
			&\operatorname{Hom}_{k[G]} \left(\operatorname{Ind}_K^G U,\operatorname{Ind}_H^G W\right)\\
			&\qquad\cong \bigoplus_{x\in X}\operatorname{Hom}_{k[L_x]} \bigl(\operatorname{Res}_{L_x}^K U, \operatorname{Res}_{L_x}^{xHx^{-1}}({}^xW)\bigr).
		\end{aligned}
	\end{equation}
\end{theorem}
\begin{proof}
	\cref{b5:thm:frobenius-reciprocity}把左侧识别为 \(\operatorname{Hom}_{k[K]}(U,\operatorname{Res}_K^G\operatorname{Ind}_H^G W)\)。 将\cref{b5:thm:mackey-decomposition}代入第二个变量，便得到各个
	\[
		\operatorname{Hom}_{k[K]} \bigl(U,\operatorname{Ind}_{L_x}^K \operatorname{Res}_{L_x}^{xHx^{-1}}({}^xW)\bigr)
	\]
	的直和。

	此时诱导出现在 Hom 的第二个变量，不能再次直接使用同一方向的左伴随。\cref{b5:thm:induction-coinduction}说明有限群的诱导与余诱导自然同构；\cref{b5:prop:coinduction-adjunction}给出的右伴随性质把这一项识别为
	\[
		\operatorname{Hom}_{k[L_x]} \bigl(\operatorname{Res}_{L_x}^K U, \operatorname{Res}_{L_x}^{xHx^{-1}}({}^xW)\bigr).
	\]
\end{proof}

在复数域上，取特征标后，公式也计算两个诱导特征标的内积。取 \(K=H\)、\(U=W=\mathbf1_H\)，每一项都是一维空间，故
\[
	\dim_{\mathbb C}\operatorname{End}_{\mathbb C[G]} \bigl(\mathbb C[G/H]\bigr)=|H\backslash G/H|.
\]
这个等式中每个双陪集对应的算子也可以直接写出。对双陪集 \(D=HxH\)，在陪集基 \(e_{gH}\) 上定义
\[
	T_D(e_{gH})=\sum_{yH\subseteq gD}e_{yH}.
\]
\(D\) 对右乘 \(H\) 稳定，所以 \(gD\) 是若干左 \(H\)-陪集之并，求和有意义；\(D\) 对左乘 \(H\) 稳定，所以 \((gh)D=gD\)，公式与 \(gH\) 所选代表元无关。对 \(a\in G\)，左乘把 \(gD\) 中的陪集逐一送到 \(agD\) 中的陪集，故 \(T_D(ae_{gH})=aT_D(e_{gH})\)。不同 \(D\) 的算子作用于 \(e_H\) 后得到支集互不相交的非零向量，因此这些算子线性无关；由上述维数等式，它们组成自同态空间的一组基。它的矩阵正是 \(G/H\times G/H\) 上轨道 \(\{(yH,gH):yH\subseteq gD\}\) 的示性矩阵，与\cref{b5:prop:permutation-character}中的轨道矩阵一致。

把 \(e_{gH}\) 识别为 \(G/H\) 上的点示性函数，就可用函数卷积记录这些算子的乘法。设 \(f,g:G\to\mathbb C\) 都在左右乘 \(H\) 下不变，\(\varphi:G/H\to\mathbb C\)，规定
\begin{equation}\label{b5:eq:finite-hecke-convolution-action}
	\begin{aligned}
		(f*g)(z)&=\frac1{|H|}\sum_{y\in G}f(y)g(y^{-1}z),\\
		(A_f\varphi)(xH)&=\frac1{|H|}\sum_{y\in G}f(x^{-1}y)\varphi(yH).
	\end{aligned}
\end{equation}
双边不变性保证第一式仍为双边不变函数，并使第二式与 \(xH\) 的代表元无关。左平移作用为 \((a\varphi)(xH)=\varphi(a^{-1}xH)\)，在求和中把 \(y\) 改为 \(ay\)，便知它与 \(A_f\) 交换。两个算子的复合满足
\[
	\begin{aligned}
		(A_fA_g\varphi)(xH) &=\frac1{|H|^2}\sum_{y,z\in G}f(x^{-1}y)g(y^{-1}z)\varphi(zH)\\
		&=\frac1{|H|}\sum_{z\in G}(f*g)(x^{-1}z)\varphi(zH),
	\end{aligned}
\]
故 \(A_fA_g=A_{f*g}\)，即卷积的次序与算子复合的次序相同。单位函数是 \(1_H\)，因为 \(A_{1_H}=I\)。按上述陪集基算子的定义，\(A_{1_D}=T_{D^{-1}}\)：其作用于 \(e_{gH}\) 后，在 \(xH\) 处非零的条件是 \(x^{-1}g\in D\)，亦即 \(xH\subseteq gD^{-1}\)。双陪集示性函数给出全部双边不变函数的一组基，而这些算子给出自同态空间的一组基，因此这个卷积代数与 \(\operatorname{End}_G\mathbb C[G/H]\) 相互识别。

\ProseParagraphBreak

在 \(H=\langle(12)\rangle\) 时的 \(S_3/H\) 三点模型中，\(T_H=I\)，而另一个双陪集的算子把每个基向量送到其余两个基向量之和。若 \(J\) 表示所有矩阵元均为 \(1\) 的三阶矩阵，则后一算子就是 \(T=J-I\)，从而自同态空间为 \(\mathbb CI\oplus\mathbb CT\)。两个双陪集均与其逆相同，所以在函数模型中 \(A_{1_H}=I\)、\(A_{1_D}=T\)。由 \(J^2=3J\)，有
\[
	T^2=(J-I)^2=J+I=T+2I, \qquad 1_D*1_D=1_D+2\cdot1_H.
\]
这里 \(|H|^{-1}\) 的归一化使 \(1_H\) 成为单位，也使非平凡双陪集的示性函数恰对应 \(J-I\)，不需要再乘 \(|H|\)。双陪集的计数在这里变成了两类具体的交织算子及其乘法。

\ProseParagraphBreak

例如，\(S_n\) 在 \(n\) 个点上的置换表示对应 \(H=S_{n-1}\)；当 \(n\geq2\) 时只有两个双陪集，所以它的自同态空间维数为 \(2\)。平凡表示出现一次，其余直和项的重数平方和因而为 \(1\)，得到一个不可约的 \((n-1)\) 维分量。

\subsection{不可约性判据}
\label{b5:subsec:0402:03}

\begin{theorem}[Mackey 不可约性判据]\label{b5:thm:mackey-irreducibility}
	设 \(G\) 为有限群，\(H\leq G\)，\(W\) 为非零有限维复 \(H\)-表示。则 \(\operatorname{Ind}_H^G W\) 不可约，当且仅当 \(W\) 不可约，且对每个不等于 \(H\) 的双陪集 \(HxH\)，令 \(L_x=H\cap xHx^{-1}\) 后有
	\[
		\operatorname{Hom}_{\mathbb C[L_x]} \bigl(\operatorname{Res}_{L_x}^H W, \operatorname{Res}_{L_x}^{xHx^{-1}}({}^xW)\bigr)=0.
	\]
\end{theorem}
\begin{proof}
	\cref{b5:thm:maschke}说明有限群的复表示完全可约。若一个非零表示分解为互不同构不可约表示 \(V_i\) 的 \(\bigoplus_iV_i^{\oplus m_i}\)，则\cref{b5:prop:isotypic-decomposition}给出其自同态空间维数为 \(\sum_i m_i^2\)；这个数等于 \(1\)，恰好意味着表示不可约。

	对\cref{b5:thm:mackey-intertwining}取 \(K=H\)、\(U=W\)，代表元 \(x=1\) 的项就是 \(\operatorname{End}_{\mathbb C[H]}W\)。所有直和项的维数非负，且这一项至少为 \(1\)。因此总维数为 \(1\)，当且仅当这一项维数为 \(1\) 而其余项都为零。再次使用完全可约性与上述维数公式，前一条件等价于 \(W\) 不可约。
\end{proof}

代表元 \(1\) 的项记录 \(W\) 原有的自同态，其余双陪集则检查 \(W\) 与共轭表示在交集上的限制之间是否还存在非零交织。只要存在这样一项，诱导表示就有额外的自同态；对有限群的复表示，完全可约性和上述重数公式说明，自同态空间维数超过 \(1\) 便迫使表示可约。例如取 \(G=S_3\)、\(H=\langle(12)\rangle\)，对非平凡双陪集仍有 \(L_{(13)}=\{1\}\)。若 \(W\) 是任意非零有限维复 \(H\)-表示，令 \(d=\dim_{\mathbb C}W\)，则这一项是两个 \(d\) 维空间之间的全部线性映射，维数为 \(d^2\)。因此
\[
	\dim_{\mathbb C}\operatorname{End}_{\mathbb C[S_3]} \bigl(\operatorname{Ind}_H^{S_3}W\bigr) \geq 1+d^2>1,
\]
所以从这个非正规子群诱导的非零表示都可约。对于前面的符号表示，代表元 \(x=(13)\) 对应的交织项维数为 \(1\)，加上单位双陪集的一维项，自同态空间的维数为 \(2\)。不可约重数的平方和因此为 \(2\)，这与\cref{b5:exm:s3-transposition-induction}中 \(\operatorname{Ind}_H^{S_3}W\cong\varepsilon\oplus U\) 中两个互不同构且各出现一次的不可约分量一致。

\ProseParagraphBreak

不可约性判据中的基域条件不能从前面的模同构公式中自动删去。在模特征下，表示可能不是完全可约的；此时自同态空间一维并不足以推出不可约。

\begin{example}\label{b5:exm:dihedral-induction}
	设 \(G=D_{2n}=\langle r,s:r^n=s^2=1,\ srs=r^{-1}\rangle\)，其中 \(n\geq3\)，且 \(H=\langle r\rangle\)。令 \(\zeta\) 为本原 \(n\) 次单位根，取 \(\lambda_j(r)=\zeta^j\)。因为 \(H\) 正规，双陪集代表元为 \(1,s\)，交集均为 \(H\)，并且 \({}^s\lambda_j=\lambda_{-j}\)。所以 \(\operatorname{Ind}_H^G\lambda_j\) 不可约，当且仅当 \(2j\not\equiv0\pmod n\)。

	\ProseParagraphBreak
	在基 \(1\otimes1,s\otimes1\) 下，两个生成元的矩阵为
	\[
		r\longmapsto
		\begin{pmatrix}\zeta^j&0\\0&\zeta^{-j}\end{pmatrix},
		\qquad s\longmapsto
		\begin{pmatrix}0&1\\1&0\end{pmatrix}.
	\]
	当 \(2j\equiv0\) 时，向量 \((1,1)\) 和 \((1,-1)\) 各自生成不变直线，诱导表示分成两个一维表示。这也从矩阵上解释了判据中的例外。
\end{example}

% !TeX root = ../../Book5.tex
% 纲要标题：### 4.3 正规子群与 Clifford 理论
% 纲要位置：工作记录/计划纲要.md，第 3831 行。
% 纲要细目（写作范围）：
% * 不可约表示限制到正规子群
% * 共轭轨道与惯性群
% * Clifford 对应


\section{正规子群与 Clifford 理论}
\label{b5:sec:0403}

一个不可约表示限制到子群后通常会分解。子群正规时，共轭作用把这些分量联系起来：它们不能来自彼此无关的若干共轭轨道。本节在复数域上证明这一结构，并由其中一个等型分量恢复整个表示。

\subsection{不可约表示限制到正规子群}
\label{b5:subsec:0403:01}

先回顾\cref{b5:prop:isotypic-decomposition}中的等型分量。对有限群 \(N\) 的有限维复表示 \(V\)，以及不可约复 \(N\)-表示 \(W\)，令 \(V_W\) 为 \(V\) 中所有同构于 \(W\) 的 \(N\)-子模之和；若没有这种子模，就令 \(V_W=0\)。

\ProseParagraphBreak

由\cref{b5:thm:maschke}，\(V\) 是各个非零 \(V_W\) 的直和，而且 \(V_W\cong W^{\oplus e_W}\)。这个子空间不依赖所选的不可约直和分解：任意简单子模到不同类型分量的投影均为零，因为\cref{b5:lem:schur}说明不同不可约表示之间没有非零同态。同样的理由说明，任意 \(N\)-子模 \(U\subseteq V\) 都有
\begin{equation}\label{b5:eq:isotypic-submodule}
U=\bigoplus_{[W]}(U\cap V_W).
\end{equation}
事实上，先把 \(U\) 分解为简单子模，每个简单子模都只落在一种等型分量中。

\ProseParagraphBreak
等式\eqref{b5:eq:isotypic-submodule}并非任意线性直和分解都有的性质。将\cref{b5:thm:central-character-idempotents}用于有限群 \(N\)，令 \(p_W\in\mathbb C[N]\) 为类型 \(W\) 对应的中心幂等元，则 \(V_W=p_WV\)。因为 \(U\) 被每个 \(n\in N\) 保持，也被这些作用的任意复线性组合保持，所以 \(p_WU\subseteq U\cap V_W\)；而 \(p_W\) 在 \(V_W\) 上为恒等作用，反向包含也成立，故 \(p_WU=U\cap V_W\)。这些投影来自群代数内部，会同时保持每个子表示；任意选取的直和投影未必保持所考察的子空间。

\begin{theorem}[Clifford 限制定理]\label{b5:thm:clifford-restriction}
设 \(G\) 为有限群，\(N\trianglelefteq G\)，\(V\) 为不可约有限维复 \(G\)-表示，\(W\) 为 \(\operatorname{Res}_N^G V\) 的一个不可约成分。则限制表示中的不可约类型恰好组成 \(W\) 在 \(G\) 共轭作用下的一个轨道，而且所有类型的重数相等。
\end{theorem}
这一结论及其后续对应构成 \term[clifford-lilun]{Clifford 理论}[Clifford theory]的基本部分。
\begin{proof}
由于 \(N\) 正规，\({}^gW\) 仍是 \(N\)-表示。若 \(A\subseteq V\) 是同构于 \(W\) 的 \(N\)-子模，则
\[
n(ga)=g(g^{-1}ng)a\qquad(n\in N,\ a\in A)
\]
说明 \(gA\cong{}^gW\)。对所有这样的 \(A\) 求和，得到
\begin{equation}\label{b5:eq:clifford-component-transport}
gV_W=V_{{}^gW}.
\end{equation}
因此，同一共轭轨道对应的等型分量之和是非零 \(G\)-子模。\(V\) 不可约，迫使这个和等于 \(V\)，故不可能出现第二个轨道。

若 \(V_W\cong W^{\oplus e}\)，则等式右边由左边的线性同构得到
\(V_{{}^gW}\cong({}^gW)^{\oplus e}\)。所以各个类型出现相同的重数 \(e\)。
\end{proof}

证明中的完全可约性用于限制到 \(N\) 的表示和它的子模，由复数域上的 Maschke 定理保证。若要推广到其他特征，需要另行检验这一步。

\subsection{共轭轨道与惯性群}
\label{b5:subsec:0403:02}

\begin{definition}\label{b5:def:inertia-group}
设 \(G\) 为有限群，\(N\trianglelefteq G\)，\(W\) 为不可约有限维复 \(N\)-表示。子集
\[
I_G(W)=\{g\in G:{}^gW\cong W\}
\]
称为 \(W\) 在 \(G\) 中的\term[guanxing-qun]{惯性群}[inertia group]。
\end{definition}
\symidx{\(I_G(W)\)：不可约表示的惯性群}

共轭构成 \(G\) 在不可约同构类上的作用，所以 \(I_G(W)\) 是子群。每个 \(n\in N\) 都属于它：若 \(W\) 的作用为 \(\rho\)，则 \(\rho(n)\) 把 \({}^nW\) 交织到 \(W\)。因此 \(N\subseteq I_G(W)\)，轨道由左陪集 \(G/I_G(W)\) 标记。

\begin{corollary}\label{b5:cor:clifford-character}
设 \(G\) 为有限群，\(N\trianglelefteq G\)，\(V\) 为不可约有限维复 \(G\)-表示，\(W\) 为 \(\operatorname{Res}_N^G V\) 的一个不可约成分。令 \(I=I_G(W)\)，令 \(e\) 为 \(W\) 在该限制表示中的重数。若 \(\chi,\theta\) 分别是 \(V,W\) 的特征标，则
\[
\operatorname{Res}_N^G V\cong
\bigoplus_{gI\in G/I}({}^gW)^{\oplus e},
\qquad
\chi|_N=e\sum_{gI\in G/I}\theta^g,
\]
其中 \(\theta^g(n)=\theta(g^{-1}ng)\)、\(n\in N\)，各和取 \(G/I\) 的一组左陪集代表元。特别地，
\[
\dim_{\mathbb C}V=e[G:I]\dim_{\mathbb C}W.
\]
\end{corollary}
\begin{proof}
\cref{b5:thm:clifford-restriction}给出一个共轭轨道及相同重数。轨道中的类型恰由 \(G/I\) 标记，因而得到模分解。对作用矩阵取迹得到特征标等式，在单位元处取值得到维数等式。
\end{proof}

例如，取 \(G=S_3\)、\(N=A_3=\langle r\rangle\)，其中 \(r=(123)\)、\(s=(12)\)。在二维不可约复表示 \(V\) 中，取 \(r\) 的特征向量 \(v_\zeta\)，使 \(rv_\zeta=\zeta v_\zeta\)，其中 \(\zeta=\exp(2\pi\mathrm i/3)\)。令 \(W=\mathbb Cv_\zeta\)。由 \(srs=r^{-1}\)，有
\[
r(sv_\zeta)=sr^{-1}v_\zeta=\zeta^{-1}sv_\zeta.
\]
两特征值不同，所以 \(V=W\oplus sW\)，而 \(V_W=W\)。换位把 \(W\) 的特征标 \(\theta\) 变成不同的 \(\theta^{-1}\)，故 \(I_{S_3}(W)=A_3\)。这正是\cref{b5:exm:dihedral-induction}中 \(n=3\) 的情形。

\subsection{Clifford 对应}
\label{b5:subsec:0403:03}

从 \(N\) 上的一种不可约类型 \(W\) 出发，直接诱导到 \(G\) 未必能恢复某一个不可约表示。例如在\cref{b5:tab:d8-q8}的四元数群中，取 \(G=Q_8\)、\(N=\{\pm1\}\)，令一维复表示 \(W\) 满足 \(-1\) 作用为 \(-1\)。\(N\) 是中心子群，故 \(I_G(W)=G\)。已构造的二维不可约复表示 \(U\) 中，\(-1\) 作用为 \(-I_U\)，所以 \(U|_N\cong W\oplus W\)。由\cref{b5:thm:frobenius-reciprocity}，
\[
\dim_{\mathbb C}\operatorname{Hom}_{\mathbb C[G]}
\bigl(\operatorname{Ind}_N^G W,U\bigr)
=\dim_{\mathbb C}\operatorname{Hom}_{\mathbb C[N]}(W,U|_N)=2.
\]
诱导表示完全可约，因而含有两份 \(U\)；它的维数为 \([G:N]\dim W=4\)，恰好得到 \(\operatorname{Ind}_N^G W\cong U\oplus U\)。

\ProseParagraphBreak
这里整个 \(W\)-等型分量是 \(U_W=U\)，它带有惯性群 \(I=G\) 的二维作用；仅保留其中一份一维 \(N\)-表示 \(W\)，就丢失了这些同型副本之间的群作用。保留 \(U_W\) 后，\(\operatorname{Ind}_I^G U_W=\operatorname{Ind}_G^G U\cong U\)。下面的对应将这种从等型分量恢复表示的构造推广到任意惯性群。

\begin{theorem}[Clifford 对应]\label{b5:thm:clifford-correspondence}
设 \(G\) 为有限群，\(N\trianglelefteq G\)，\(W\) 为不可约有限维复 \(N\)-表示，\(I=I_G(W)\)。诱导给出以下两类有限维表示的同构类之间的双射：
\begin{enumerate}
\item 限制到 \(N\) 时含有 \(W\) 的不可约复 \(I\)-表示 \(U\)；
\item 限制到 \(N\) 时含有 \(W\) 的不可约复 \(G\)-表示 \(V\)。
\end{enumerate}
正向映射为 \(U\mapsto\operatorname{Ind}_I^G U\)，逆映射为 \(V\mapsto V_W\)，其中 \(V_W\) 为 \(\operatorname{Res}_N^G V\) 中的 \(W\)-等型分量。特别地，\(V_W\) 是不可约 \(I\)-表示，而自然映射
\begin{equation}\label{b5:eq:clifford-reconstruction}
\mathbb C[G]\otimes_{\mathbb C[I]}V_W\longrightarrow V,
\qquad g\otimes v\longmapsto gv
\end{equation}
是 \(G\)-同构。
\end{theorem}
\begin{proof}
先给定第二类中的 \(V\)。等式\eqref{b5:eq:clifford-component-transport}说明 \(V_W\) 是 \(I\)-不变的。若 \(0\neq U\subseteq V_W\) 是 \(I\)-子模，则选取 \(G/I\) 的代表元后，\(\bigoplus_{gI}gU\) 是 \(G\)-子模：任何群元素左乘代表元，只会改变代表元并带出一个 \(I\) 中的因子。这个和非零，所以等于 \(V\)。不同 \(gU\) 落在不同的 \(N\)-等型分量，取其中的 \(W\)-等型部分便有 \(U=V_W\)。故 \(V_W\) 不可约。

映射\eqref{b5:eq:clifford-reconstruction}的平衡关系由 \(I\)-作用保证。它的像包含所有 \(gV_W\)，故满射；源端的各个陪集分量映到目标中不同的等型分量，而且每一分量上的映射均为线性同构。因此它也单射。

反过来，设 \(U\) 属于第一类。对有限群 \(I\) 及其正规子群 \(N\) 应用\cref{b5:thm:clifford-restriction}，由于 \(I\) 固定 \(W\) 的同构类，得到 \(U|_N\cong W^{\oplus e}\)。令 \(Y=\operatorname{Ind}_I^G U\)。作为 \(N\)-模，它的陪集分解为
\[
Y=\bigoplus_{gI\in G/I}g\otimes U,
\]
每一项的类型是 \({}^gW\)，不同陪集的类型互不相同。

若 \(0\neq Z\subseteq Y\) 是 \(G\)-子模，等式\eqref{b5:eq:isotypic-submodule}说明它与某个 \(g\otimes U\) 有非零交。用 \(g^{-1}\) 作用，把非零交移到 \(1\otimes U\)。于是 \(Z\cap(1\otimes U)\) 是非零 \(I\)-子模，必等于 \(1\otimes U\)。再用 \(G\) 作用，得到 \(Z=Y\)。所以 \(Y\) 不可约。它的 \(W\)-等型分量恰为 \(1\otimes U\)，故先诱导再取 \(W\)-等型分量恢复 \(U\)；反方向已由 \(\operatorname{Ind}_I^G V_W\simeq V\) 得到。若两个 \(I\)-模同构，诱导后仍同构；若两个 \(G\)-模同构，其同构把 \(W\)-等型分量对应起来。因此两种构造在同构类上良定义且互为逆。
\end{proof}

回到上面的 \(S_3\) 例子，此时 \(I=N\)、\(V_W=W=\mathbb Cv_\zeta\)。左陪集代表取 \(1,s\)，则
\[
\operatorname{Ind}_N^{S_3}W
=\mathbb C(1\otimes v_\zeta)\oplus\mathbb C(s\otimes v_\zeta).
\]
映射\eqref{b5:eq:clifford-reconstruction}将这组基分别送到 \(v_\zeta,sv_\zeta\)，后者正是 \(V=W\oplus sW\) 的一组基。在这两组对应的基下，生成元的作用矩阵相同：
\[
r\longmapsto\begin{pmatrix}\zeta&0\\0&\zeta^{-1}\end{pmatrix},
\qquad
s\longmapsto\begin{pmatrix}0&1\\1&0\end{pmatrix}.
\]
其中 \(r\) 保持两个陪集分量，\(s\) 交换它们；它们在 \(V\) 中对应的就是两条不同类型的 \(N\)-特征直线。诱导的陪集分解在这里逐项重建了原表示。

\ProseParagraphBreak
Clifford 对应把问题从 \(G\) 缩到惯性群 \(I\)。当 \(I=N\) 时，\(W\) 本身已经诱导出不可约 \(G\)-表示；当 \(I=G\) 时，限制只含一种类型，但重数 \(e\) 仍可能大于 \(1\)。后一种情形需要研究 \(W\) 能否延拓，以及重数空间上究竟有怎样的作用。

% !TeX root = ../../Book5.tex
% 纲要标题：### 4.4 射影表示初步
% 纲要位置：工作记录/计划纲要.md，第 4156 行。
% 纲要细目（写作范围）：
% * 延拓表示的障碍
% * 障碍余圈、射影表示的因子集与 2-余圈
% * 与第四卷群上同调和第二卷交叉积的联系
% * 交换群的扭曲乘法、中心及矩阵代数算例；障碍类对重数的维数限制
% ---


\section{射影表示初步}
\label{b5:sec:0404}

惯性群中的元素都把 \(W\) 变成同构的表示，但为每个元素分别选一个同构，并不保证这些同构按照群乘法相乘。差出的标量恰好记下延拓的障碍。本节从这些标量出发研究这一问题。

\subsection{表示延拓的障碍}
\label{b5:subsec:0404:01}

\begin{definition}\label{b5:def:representation-extension}
设 \(I\) 为有限群，\(N\trianglelefteq I\)，\(W\) 为有限维复向量空间，\(\rho:N\to\operatorname{GL}_{\mathbb C}(W)\) 为表示。若表示 \(\widetilde\rho:I\to\operatorname{GL}_{\mathbb C}(W)\) 在 \(N\) 上等于 \(\rho\)，则称它为 \(\rho\) 到 \(I\) 的\term[biaoshi-yantuo]{延拓}[extension of a representation]。
\end{definition}

延拓保持原来的向量空间，维数不变；它与一般会增大维数的诱导是不同的构造。如果 \(\rho\) 能延拓，则对 \(g\in I\)，矩阵 \(\widetilde\rho(g)\) 实现 \(W\cong{}^gW\)，故 \(W\) 的同构类必须被 \(I\) 固定。下面精确考察这个必要条件还差什么。

\begin{theorem}[延拓的标量障碍]\label{b5:thm:extension-obstruction}
设 \(I\) 为有限群，\(N\trianglelefteq I\)，\(\rho:N\to\operatorname{GL}_{\mathbb C}(W)\) 为不可约有限维复表示，并假定 \({}^gW\cong W\) 对所有 \(g\in I\) 成立。令 \(Q=I/N\)，选取商映射的集合截面 \(s:Q\to I\)，满足 \(s(1)=1\)，并对 \(p,q\in Q\) 记
\[
c(p,q)=s(p)s(q)s(pq)^{-1}\in N.
\]
可以选取 \(T_q\in\operatorname{GL}_{\mathbb C}(W)\)，使 \(T_1=1_W\)，且
\begin{equation}\label{b5:eq:extension-intertwiner}
T_q\rho(n)T_q^{-1}
=\rho\bigl(s(q)ns(q)^{-1}\bigr)\qquad(n\in N).
\end{equation}
存在唯一标量 \(\alpha(p,q)\in\mathbb C^\times\)，使
\begin{equation}\label{b5:eq:extension-factor}
T_pT_q=\alpha(p,q)\rho\bigl(c(p,q)\bigr)T_{pq}.
\end{equation}
这些标量满足
\begin{equation}\label{b5:eq:factor-cocycle}
\alpha(1,q)=\alpha(q,1)=1,\qquad
\alpha(p,q)\alpha(pq,r)=\alpha(q,r)\alpha(p,qr)\quad(p,q,r\in Q).
\end{equation}
称 \(\alpha\) 为这组选择所确定的障碍余圈。
表示 \(W\) 能延拓到 \(I\)，当且仅当存在 \(b:Q\to\mathbb C^\times\)、\(b(1)=1\)，使
\begin{equation}\label{b5:eq:extension-rescaling}
b(p)b(q)\alpha(p,q)=b(pq)\qquad(p,q\in Q).
\end{equation}
\end{theorem}
\begin{proof}
先比较截面元素的两种三项乘积。将 \((s(p)s(q))s(r)\) 与 \(s(p)(s(q)s(r))\) 都写成 \(N\) 中元素乘 \(s(pqr)\)，消去末尾的 \(s(pqr)\)，便得到
\[
 c(p,q)c(pq,r)
 =s(p)c(q,r)s(p)^{-1}c(p,qr).
\]
所以 \(c\) 满足的是带共轭修正的恒等式：它记录截面不保持乘法的偏差，取值于可能非交换的 \(N\)。下面的 \(\alpha\) 则是在用 \(\rho(c(p,q))\) 修正这一偏差以后，交织算子乘法中剩余的标量；两者不能按同一条标量余圈公式处理。

共轭不变性给出满足\eqref{b5:eq:extension-intertwiner}的可逆交织映射。\(T_pT_q\) 和 \(\rho\bigl(c(p,q)\bigr)T_{pq}\) 都实现由 \(s(p)s(q)\) 所给的共轭，所以二者之比与 \(\rho(N)\) 对易。\cref{b5:cor:schur-algebraically-closed}说明该比值为唯一非零标量，得到\eqref{b5:eq:extension-factor}。

每个 \(g\in I\) 唯一写成 \(ns(q)\)。定义可逆算子
\[
P\bigl(ns(q)\bigr)=\rho(n)T_q.
\]
若 \(g=ns(p)\)、\(h=ms(q)\)，则
\[
\begin{aligned}
P(g)P(h)
&=\rho\bigl(n s(p)m s(p)^{-1}\bigr)T_pT_q\\
&=\alpha(p,q)
\rho\bigl(n s(p)m s(p)^{-1}c(p,q)\bigr)T_{pq}\\
&=\alpha(p,q)P(gh).
\end{aligned}
\]
对三个元素比较 \(\bigl(P(g)P(h)\bigr)P(\ell)\) 与 \(P(g)\bigl(P(h)P(\ell)\bigr)\)，并消去可逆的 \(P(gh\ell)\)，就得到\eqref{b5:eq:factor-cocycle}中的第二个等式；单位元处的等式来自 \(T_1=1_W\)。

如果存在 \(b\) 满足\eqref{b5:eq:extension-rescaling}，则
\(\widetilde\rho\bigl(ns(q)\bigr)=\rho(n)b(q)T_q\) 按上述乘法计算是群表示，而且在 \(N\) 上等于 \(\rho\)。反过来，若有延拓，则 \(\widetilde\rho\bigl(s(q)\bigr)\) 与 \(T_q\) 实现同一个共轭。Schur 引理给出 \(\widetilde\rho\bigl(s(q)\bigr)=b(q)T_q\)，再比较两者乘法便得到\eqref{b5:eq:extension-rescaling}。
\end{proof}

不能把“不变”直接换成“可延拓”。取
\[
I=Q_8=\{\pm1,\pm i,\pm j,\pm ij\},\qquad
N=\{\pm1\},
\]
并令 \(\rho(-1)=-1\)。\(N\) 在中心，故这个一维表示被全部 \(I\) 固定。但任何一维 \(Q_8\)-表示都把交换子送到 \(1\)，而 \(iji^{-1}j^{-1}=-1\)，所以它不能延拓。这个障碍在下一小节可以直接算成一个不对称的标量函数。

\subsection{因子集与 2-余圈}
\label{b5:subsec:0404:02}

\begin{definition}\label{b5:def:projective-representation}
设 \(Q\) 为有限群，\(E\neq0\) 为有限维复向量空间。若映射 \(R:Q\to\operatorname{GL}_{\mathbb C}(E)\) 满足 \(R(1)=1_E\)，并对每个 \(p,q\in Q\) 存在非零标量 \(\alpha(p,q)\)，使
\[
R(p)R(q)=\alpha(p,q)R(pq),
\]
则称 \(R\) 为 \(Q\) 的\term[sheying-biaoshi]{射影表示}[projective representation]，称 \(\alpha\) 为所选矩阵的\term[yinziji]{因子集}[factor set]。若 \(E\) 没有非零真子空间同时被所有 \(R(q)\) 保持，则称这个射影表示不可约。
对于因子集相同的两个射影表示 \(R:Q\to\operatorname{GL}_{\mathbb C}(E)\)、\(R':Q\to\operatorname{GL}_{\mathbb C}(E')\)，若存在线性同构 \(a:E\to E'\)，使 \(aR(q)=R'(q)a\) 对所有 \(q\in Q\) 成立，则称它们同构。
\end{definition}

把矩阵模去非零标量后，射影表示给出群同态
\(Q\to\operatorname{PGL}(E)\)。反之，这样的群同态选取矩阵代表后就给出因子集；代表不同，因子集也可能不同。

\ProseParagraphBreak
这里模去标量的作用，比定义中的同构关系更粗。例如 \(Q=C_2=\{1,s\}\) 的两种一维表示 \(R_+(s)=1\)、\(R_-(s)=-1\) 都有平凡因子集，而且 \(\operatorname{PGL}_1(\mathbb C)=\{1\}\)，所以给出相同的射影作用。但若有线性同构 \(a\) 交织它们，就须 \(a=-a\)，与 \(a\ne0\) 矛盾。固定因子集的同构要求逐个矩阵严格交织；在 \(\operatorname{PGL}\) 中相同或共轭的作用则允许逐个改变标量。在这个例子中，\(R_-\) 是用非平凡一维特征标扭曲 \(R_+\) 得到的；扭曲没有改变平凡因子集，却使它们成为不同构的表示。因此因子集相同并不蕴含射影表示同构。

\begin{definition}\label{b5:def:scalar-two-cocycle}
设 \(Q\) 为有限群，\(\alpha:Q\times Q\to\mathbb C^\times\) 为函数。若对所有 \(p,q,r\in Q\) 有
\[
\alpha(1,q)=\alpha(q,1)=1,\qquad
\alpha(p,q)\alpha(pq,r)=\alpha(q,r)\alpha(p,qr),
\]
则称 \(\alpha\) 为系数作用平凡的\term[guiyihua-eryuquan]{归一化 \(2\)-余圈}[normalized 2-cocycle]。若存在 \(a:Q\to\mathbb C^\times\)、\(a(1)=1\)，使
\[
\alpha(p,q)=\frac{a(p)a(q)}{a(pq)}\qquad(p,q\in Q),
\]
则称 \(\alpha\) 为一个\term[eryubian]{\(2\)-余边}[2-coboundary]。
\end{definition}

矩阵乘法的结合律说明每个因子集都是归一化 \(2\)-余圈。将 \(R(q)\) 改成 \(b(q)R(q)\)，其中 \(b:Q\to\mathbb C^\times\)、\(b(1)=1\)，因子集变成
\begin{equation}\label{b5:eq:factor-coboundary-change}
\alpha'(p,q)=\frac{b(p)b(q)}{b(pq)}\alpha(p,q).
\end{equation}
在\cref{b5:thm:extension-obstruction}中，\(T_pT_q\) 与 \(T_{pq}\) 之间还差一个算子 \(\rho(c(p,q))\)，它一般不是标量。因此 \(q\mapsto T_q\) 一般不是 \(Q\) 的射影表示。证明中构造的
\[
P\bigl(ns(q)\bigr)=\rho(n)T_q
\]
才是 \(I\) 上的射影表示，它满足
\[
P(g)P(h)=\alpha(gN,hN)P(gh)\qquad(g,h\in I).
\]
其因子集是把 \(Q\) 上的障碍余圈 \(\alpha\) 沿商映射拉回到 \(I\) 得到的函数；这个拉回称为余圈的膨胀（inflation）。延拓存在的条件则是 \(\alpha\) 本身为一个 \(2\)-余边。

\ProseParagraphBreak
非标量因子可以在 \(I=S_3\times C_2\)、\(N=S_3\times\{1\}\) 中直接看见。令 \(W\) 为 \(S_3\) 的二维标准表示 \(\rho\)，记 \(C_2\) 的非平凡元为 \(q\)，并取 \(r=(123)\)、\(s(q)=(r,q)\)。选 \(T_q=\rho(r)\)、\(T_1=I_W\)。这些算子实现截面元素的共轭，而且 \(c(q,q)=(r^2,1)\)，所以
\[
T_q^2=\rho(r^2),\qquad \alpha(q,q)=1.
\]
在一组特征向量基下，\(\rho(r)=\operatorname{diag}(\omega,\omega^2)\)，其中 \(\omega=\exp(2\pi\mathrm i/3)\)，故 \(\rho(r^2)=\operatorname{diag}(\omega^2,\omega)\) 不是标量。这里 \(W\) 实际上可延拓到直积群，障碍余圈处处为一，但这个截面给出的 \(q\mapsto T_q\) 仍不是 \(C_2\) 的射影表示：\(T_q^2\) 不是 \(T_1\) 的标量倍数。截面乘法中的 \(N\)-因子与标量障碍因而必须分别记录。

\begin{proposition}\label{b5:prop:factor-choice-independence}
设 \(I\) 为有限群，\(N\trianglelefteq I\)，\(\rho:N\to\operatorname{GL}_{\mathbb C}(W)\) 为同构类被 \(I\) 固定的不可约有限维复表示。令 \(Q=I/N\)，取商映射的截面 \(s:Q\to I\)、\(s(1)=1\)，以及算子 \(T_q\in\operatorname{GL}_{\mathbb C}(W)\)，使 \(T_1=1_W\) 且 \(T_q\rho(n)T_q^{-1}=\rho(s(q)ns(q)^{-1})\) 对 \(q\in Q,n\in N\) 成立。对 \(p,q\in Q\)，记 \(c(p,q)=s(p)s(q)s(pq)^{-1}\)，由 \(T_pT_q=\alpha(p,q)\rho(c(p,q))T_{pq}\) 确定障碍余圈 \(\alpha:Q\times Q\to\mathbb C^\times\)。更换交织算子或截面后，新障碍余圈 \(\alpha'\) 满足 \(\alpha'/\alpha\) 是一个 \(2\)-余边。因此，“能否延拓”的判断不依赖这些选择。
\end{proposition}
\begin{proof}
固定截面时，任意另一组交织算子由 Schur 引理写成 \(T'_q=b(q)T_q\)，其障碍余圈由\eqref{b5:eq:factor-coboundary-change}给出。

若改用 \(s'(q)=d(q)s(q)\)，其中 \(d(q)\in N\)、\(d(1)=1\)，先选 \(T'_q=\rho\bigl(d(q)\bigr)T_q\)。对同一个群元素，写成 \(ns(q)=n's'(q)\) 就有 \(n=n'd(q)\)，从而 \(\rho(n)T_q=\rho(n')T'_q\)。所以证明中定义的 \(P(g)\) 完全不变，其因子集和 \(Q\) 上的障碍余圈也不变。再作其他交织算子选择时，才按前一段改变一个余边。
\end{proof}

回到 \(Q_8\)。把 \(Q=Q_8/N\) 写成 \(\mathbb F_2^2\)，取 \(s(a,b)=i^aj^b\) 及 \(T_{(a,b)}=1\)，用 \(0,1\) 代表各个指数。计算 \(i^aj^bi^cj^d\) 时，先将 \(j^b\) 与 \(i^c\) 交换，产生 \((-1)^{bc}\)；再由 \(i^2=j^2=-1\)，将两组指数约化模二，分别产生 \((-1)^{ac}\) 与 \((-1)^{bd}\)。所以
\[
s(a,b)s(c,d)=(-1)^{ac+bd+bc}s(a+c,b+d).
\]
右边的加法在 \(\mathbb F_2^2\) 中进行。这给出 \(c\bigl((a,b),(c,d)\bigr)=(-1)^{ac+bd+bc}\)。因 \(\rho(-1)=-1\) 且 \(T_q=1\)，\eqref{b5:eq:extension-factor}要求 \(\alpha(p,q)=\rho(c(p,q))^{-1}\)，于是
\[
\alpha\bigl((a,b),(c,d)\bigr)=(-1)^{ac+bd+bc}.
\]
指数按模 \(2\) 计算。这是双线性式，所以直接展开也可验证余圈恒等式。取 \(p=(1,0)\)、\(q=(0,1)\)，得到 \(\alpha(p,q)=1\) 而 \(\alpha(q,p)=-1\)。交换群上的余边 \(a(p)a(q)/a(pq)\) 对称，故这个障碍余圈不是余边。

\subsection{与群上同调、交叉积的联系}
\label{b5:subsec:0404:03}

归一化 \(2\)-余圈在逐点乘法下组成交换群：乘积和倒数仍满足同一余圈恒等式。对函数 \(a:Q\to\mathbb C^\times\)、\(a(1)=1\)，记
\[
(\delta a)(p,q)=\frac{a(p)a(q)}{a(pq)}.
\]
它在单位元处归一化，而且
\[
(\delta a)(p,q)(\delta a)(pq,r)
=\frac{a(p)a(q)a(r)}{a(pqr)}
=(\delta a)(q,r)(\delta a)(p,qr),
\]
所以每个 \(2\)-余边都是 \(2\)-余圈。又有
\[
(\delta a)(\delta b)=\delta(ab),\qquad
(\delta a)^{-1}=\delta(a^{-1}),
\]
故全部 \(2\)-余边构成余圈群的子群。将余圈群模去这个子群，就得到记为 \(H^2(Q,\mathbb C^\times)\) 的群。这里 \(\mathbb C^\times\) 上的 \(Q\)-作用是平凡的。它是群上同调的一个具体低阶形式；当前的延拓障碍就是 \([\alpha]\)。第四卷12.1节由群余链复形定义一般的群上同调，12.2.1小节写出次数二的余圈与余边公式；将其中的加法系数改用乘法记号，再取平凡作用，就得到这里的两个公式。
% 跨卷公式位置：b4:eq:group-two-cocycle、b4:eq:group-two-coboundary（12.2.1，群的二次余圈与余边公式）。

\ProseParagraphBreak

一个余圈还可规定一个代数的乘法。给定有限群 \(Q\) 和归一化 \(2\)-余圈 \(\alpha\)，在以 \(u_q\) 为基的复向量空间上定义
\[
u_pu_q=\alpha(p,q)u_{pq}.
\]
归一化条件使 \(u_1\) 为单位元，余圈恒等式恰好使乘法结合；各个 \(u_q\) 都可逆。这个代数称为\term[niuqu-qundaisu]{扭曲群代数}[twisted group algebra]，记为 \(\mathbb C^\alpha[Q]\)。它是系数域上作用平凡的交叉积。第二卷18.4.1小节的交叉积乘法还包含 Galois 群对系数域的自同构作用；把该作用改成平凡作用，结合律便成为当前的余圈恒等式。
% 跨卷公式位置：b2:def:galois-crossed-product、b2:eq:crossed-product-multiplication（18.4.1，带 Galois 系数作用的交叉积乘法）。
给出一个非零有限维左 \(\mathbb C^\alpha[Q]\)-模，等价于给出因子集为 \(\alpha\) 的射影表示：令 \(u_q\) 按 \(R(q)\) 作用即可。重标定基 \(u_q\) 与\eqref{b5:eq:factor-coboundary-change}完全对应。
\symidx{\(\mathbb C^\alpha[Q]\)：扭曲群代数}

\ProseParagraphBreak
普通群代数在群交换时也交换，但扭曲乘法还要比较 \(\alpha(p,q)\) 与 \(\alpha(q,p)\)。这两个标量的比值决定哪些基元素能够对易，也决定中心中可以保留哪些系数。

\begin{proposition}\label{b5:prop:twisted-commutator-center}
设 \(Q\) 为有限交换群，\(\alpha:Q\times Q\to\mathbb C^\times\) 为归一化 \(2\)-余圈。记
\[
 \beta_\alpha(p,q)=\frac{\alpha(p,q)}{\alpha(q,p)}.
\]
在扭曲群代数 \(\mathbb C^\alpha[Q]\) 中，有
\[
 u_pu_q=\beta_\alpha(p,q)u_qu_p.
\]
其中心为
\[
 Z(\mathbb C^\alpha[Q])
 =\operatorname{span}_{\mathbb C}\{\,u_p:\beta_\alpha(p,q)=1\;\text{对所有}\;q\in Q\,\}.
\]
将 \(\alpha\) 改变一个 \(2\)-余边，不改变 \(\beta_\alpha\)。特别地，若 \(Q=\mathbb F_2^2\)，且
\(\alpha((a,b),(c,d))=(-1)^{bc}\)，则
\[
 Z(\mathbb C^\alpha[Q])=\mathbb C u_{(0,0)},\qquad
 \mathbb C^\alpha[Q]\cong M_2(\mathbb C).
\]
这一因子集的不可约射影表示只有一个同构类型，维数为 \(2\)。
\end{proposition}
\begin{proof}
因为 \(pq=qp\)，两个乘积的基向量相同，比较系数即得对易公式。写 \(z=\sum_p a_pu_p\)。对固定的 \(q\)，\(p\mapsto pq\) 是双射，因此 \(zu_q=u_qz\) 等价于逐个满足
\[
 a_p\bigl(\alpha(p,q)-\alpha(q,p)\bigr)=0.
\]
这对所有 \(q\) 成立，恰好给出所述中心。交换群上的余边 \((\delta b)(p,q)=b(p)b(q)/b(pq)\) 关于 \(p,q\) 对称，所以它在比值中抵消，\(\beta_\alpha\) 不变。

在给定的 \(\mathbb F_2^2\) 算例中，
\[
 \beta_\alpha((a,b),(c,d))=(-1)^{bc+ad}.
\]
分别取 \((c,d)=(1,0),(0,1)\)，要使此值总为 \(1\)，必须 \(b=a=0\)，故中心只有单位元的标量倍数。再令
\[
 X=\begin{pmatrix}1&0\\0&-1\end{pmatrix},\qquad
 Y=\begin{pmatrix}0&1\\1&0\end{pmatrix}.
\]
由于 \(X^2=Y^2=I_2\)、\(YX=-XY\)，映射 \(u_{(a,b)}\mapsto X^aY^b\) 保持乘法。四个矩阵 \(I_2,X,Y,XY\) 线性无关，张成全部二阶矩阵，所以该映射是代数同构。于是这个交换群的扭曲群代数之所以不交换，正是因为 \(\beta_\alpha\) 并非处处为 \(1\)；它的中心也与矩阵代数的标量中心一致。

最后，设 \(E\) 是非零简单 \(M_2(\mathbb C)\)-模。矩阵单位满足 \(E_{11}+E_{22}=1\)，若 \(E_{11}E\ne0\)，可直接取 \(0\ne v\in E_{11}E\)。否则由 \(E_{11}+E_{22}=1\) 取 \(0\ne w\in E_{22}E\)，令 \(v=E_{12}w\)；因为 \(E_{21}E_{12}w=w\)，有 \(v\ne0\)，而 \(E_{11}v=E_{11}E_{12}w=v\)，故仍得到 \(0\ne v\in E_{11}E\)。又 \(E_{12}E_{21}v=v\)，所以 \(v,E_{21}v\) 分别在两个不同的矩阵单位分量中，且都非零，因而线性无关。它们张成的二维子空间被全部矩阵单位保持，由简单性等于 \(E\)；在这组基下作用就是标准矩阵作用。这证明了简单模的唯一同构类型及其维数，亦即相应射影表示的结论。
\end{proof}
\symidx{\(\beta_\alpha(p,q)\)：交换群的扭曲对易标量}

\begin{theorem}[重数空间上的射影表示]\label{b5:thm:clifford-projective-multiplicity}
设 \(I\) 为有限群，\(N\trianglelefteq I\)，\(\rho:N\to\operatorname{GL}_{\mathbb C}(W)\) 为同构类被 \(I\) 固定的不可约有限维复表示。令 \(Q=I/N\)，取商映射的截面 \(s:Q\to I\)、\(s(1)=1\)，以及算子 \(T_q\in\operatorname{GL}_{\mathbb C}(W)\)，使 \(T_1=1_W\) 且
\[
T_q\rho(n)T_q^{-1}=\rho(s(q)ns(q)^{-1})\qquad(n\in N,\ q\in Q).
\]
对 \(p,q\in Q\)，记 \(c(p,q)=s(p)s(q)s(pq)^{-1}\)，以
\[
T_pT_q=\alpha(p,q)\rho(c(p,q))T_{pq}
\]
确定障碍余圈 \(\alpha:Q\times Q\to\mathbb C^\times\)。固定上述 \(s\) 与 \(T_q\)。对每个正整数 \(e\)，满足 \(\operatorname{Res}_N^I U\cong W^{\oplus e}\) 的有限维复 \(I\)-表示 \(U\) 的同构类，与 \(e\) 维复向量空间 \(E\) 上因子集为 \(\alpha^{-1}\) 的射影 \(Q\)-表示 \(S\) 的同构类一一对应。由 \(S:Q\to\operatorname{GL}_{\mathbb C}(E)\) 构造的 \(I\)-表示作用为
\begin{equation}\label{b5:eq:clifford-projective-tensor}
ns(q)\longmapsto \rho(n)T_q\otimes S(q)
\quad\hbox{作用在 }W\otimes_{\mathbb C}E\hbox{ 上}.
\end{equation}
而且 \(U\) 不可约，当且仅当 \(S\) 不可约。
\end{theorem}
\begin{proof}
先给定 \(U\)，记其表示为 \(\pi\)，并令
\(E=\operatorname{Hom}_{\mathbb C[N]}(W,U)\)。由于 \(U|_N\cong W^{\oplus e}\) 且 \(\operatorname{End}_N(W)=\mathbb C\)，求值映射
\[
W\otimes E\longrightarrow U,\qquad w\otimes f\longmapsto f(w)
\]
由\cref{b5:prop:isotypic-decomposition}是 \(N\)-同构，且 \(\dim E=e\)。要把 \(s(q)\) 的作用写成 \(T_q\otimes S(q)\)，必须让求值映射保持这一作用，也就是要求
\[
 (S(q)f)(T_qw)=\pi(s(q))f(w).
\]
由于 \(T_q\) 可逆，把 \(T_qw\) 换成任意向量，就得到
\[
S(q)f=\pi\bigl(s(q)\bigr)\,f\,T_q^{-1}.
\]
逆算子由求值的等变性所迫使，而非另作的选择。第一个张量因子的乘法产生 \(\alpha(p,q)\)，因此第二个因子必须产生其倒数，才能恢复普通的 \(I\)-作用。下面验证这些算子确实在 \(E\) 上作用，并满足这一乘法关系。
由\eqref{b5:eq:extension-intertwiner}，对 \(n\in N\) 有
\[
\begin{aligned}
\bigl(S(q)f\bigr)\rho(n)
&=\pi\bigl(s(q)\bigr)f
\rho\bigl(s(q)^{-1}ns(q)\bigr)T_q^{-1}\\
&=\pi(n)\pi\bigl(s(q)\bigr)fT_q^{-1}
=\pi(n)\bigl(S(q)f\bigr).
\end{aligned}
\]
所以此复合仍是 \(N\)-线性映射，确实属于 \(E\)。求值映射把 \(T_q\otimes S(q)\) 的作用变成 \(\pi\bigl(s(q)\bigr)\)，从而得到\eqref{b5:eq:clifford-projective-tensor}。

在求值同构 \(W\otimes E\cong U\) 下，\(\pi(s(q))=T_q\otimes S(q)\)，而 \(\pi(n)=\rho(n)\otimes1_E\)。将 \(p,q\) 处的算子相乘，群内等式 \(s(p)s(q)=c(p,q)s(pq)\) 与\eqref{b5:eq:extension-factor}给出
\[
\alpha(p,q)\rho(c(p,q))T_{pq}\otimes S(p)S(q)
=\rho(c(p,q))T_{pq}\otimes S(pq).
\]
左乘 \((\rho(c(p,q))T_{pq})^{-1}\otimes1_E\)，两边就有相同的第一因子 \(1_W\)。由于 \(W\ne0\)，映射 \(B\mapsto1_W\otimes B\) 是单射，故得到
\[
S(p)S(q)=\alpha(p,q)^{-1}S(pq).
\]
反过来，若 \(S\) 满足此等式，则\eqref{b5:eq:clifford-projective-tensor}中两个因子产生的标量互相抵消，因而给出普通的 \(I\)-表示；在 \(N\) 上它是 \(\rho(n)\otimes1_E\)。

若 \(a:U\to U'\) 为 \(I\)-表示同构，则 \(f\mapsto a\circ f\) 给出重数空间的同构，并交织两个重数空间上的射影作用。反过来，射影表示同构 \(b:E\to E'\) 给出 \(I\)-表示同构 \(1_W\otimes b\)。对于由 \(S\) 构造的 \(W\otimes E\)，识别
\[
E\xrightarrow{\sim}\operatorname{Hom}_N(W,W\otimes E),
\qquad v\longmapsto(w\longmapsto w\otimes v)
\]
把重数空间上的作用恢复为原来的 \(S\)；对于 \(U\)，上述求值同构又把构造出的 \(W\otimes E\) 送回 \(U\)。因此两种构造在同构类上互为逆。

同型副本之间的子模可以斜着嵌入，但这种混合由重数空间记录。例如 \(E=\mathbb C^2\) 时，\(W\otimes E\cong W\oplus W\)，对角子模 \(\{(w,w):w\in W\}\) 就是 \(W\otimes\mathbb C(1,1)\)，而 \(\{(w,2w):w\in W\}\) 对应 \(W\otimes\mathbb C(1,2)\)。这里的直线在 \(E\) 中，\(N\) 仍在每份 \(W\) 上按原表示作用。

一般地，设 \(Z\subseteq W\otimes E\) 为 \(N\)-子模。复表示的完全可约性说明 \(Z\) 是简单模的直和，而它在 \(W^{\oplus e}\) 中的嵌入使每个简单成分都同构于 \(W\)。通过上面的识别，将
\[
F=\operatorname{Hom}_N(W,Z)
\subseteq\operatorname{Hom}_N(W,W\otimes E)\cong E
\]
看成 \(E\) 的子空间。由 \(\operatorname{End}_N(W)=\mathbb C\)，求值映射 \(W\otimes F\to Z\) 是同构，且它与 \(Z\) 在 \(W\otimes E\) 中的包含相容，故在这个既定空间内有 \(Z=W\otimes F\)。由于每个 \(T_q\) 可逆，
\[
(T_q\otimes S(q))(W\otimes F)=W\otimes S(q)F.
\]
于是 \(W\otimes F\) 是 \(I\)-子模，当且仅当 \(F\) 被所有 \(S(q)\) 保持。这就证明了不可约性的等价。
\end{proof}

结合\cref{b5:thm:clifford-correspondence}，先在惯性群中求这些射影表示，再诱导到 \(G\)，便得到限制中含有 \(W\) 的全部不可约表示。若障碍为零，可以重标定 \(T_q\)，使它们给出 \(W\) 的延拓；此时 \(\alpha=1\)，重数空间上的表示也成为普通的商群表示。

\begin{corollary}\label{b5:cor:projective-dimension-obstruction}
设 \(Q\) 为有限群，\(R:Q\to\operatorname{GL}_{\mathbb C}(E)\) 为因子集 \(\alpha\) 的射影表示，\(d=\dim_{\mathbb C}E>0\)。则 \(\alpha^d\) 是 \(2\)-余边，因而 \([\alpha]^d=1\)；若 \([\alpha]\) 的阶为 \(m\)，则 \(m\mid d\)。
\end{corollary}
\begin{proof}
对 \(R(p)R(q)=\alpha(p,q)R(pq)\) 取行列式，得到
\[
 \alpha(p,q)^d
 =\frac{\det R(p)\det R(q)}{\det R(pq)}.
\]
右边是函数 \(q\mapsto\det R(q)\) 的余边，故 \([\alpha]^d=1\)。因此这个类有有限阶，且其阶整除 \(d\)。
\end{proof}

将推论用于\cref{b5:thm:clifford-projective-multiplicity}中的重数空间，因子集为 \(\alpha^{-1}\)，维数为 \(e\)，便有 \([\alpha]^e=1\)。所以非平凡延拓障碍会排除重数 \(e=1\)，其阶为 \(m\) 时还要求 \(m\mid e\)。但重数大于 \(1\) 不能反推出障碍非平凡。例如取 \(I=C_2\times S_3\)、\(N=C_2\times\{1\}\)，令 \(W\) 为 \(C_2\) 的非平凡一维表示，\(E\) 为 \(S_3\) 的二维标准表示。\(W\) 可通过让第二因子平凡作用延拓到 \(I\)，所以障碍为零；上述定理仍给出不可约 \(I\)-表示 \(U=W\otimes E\)，而 \(U|_N\cong W^{\oplus2}\)。实际重数还取决于商群上相应射影不可约表示的维数。

\ProseParagraphBreak
在同一个 \(Q_8\) 例子中，取二维不可约表示 \(U\)，使
\[
\pi(i)=\begin{pmatrix}\mathrm i&0\\0&-\mathrm i\end{pmatrix},
\qquad
\pi(j)=\begin{pmatrix}0&1\\-1&0\end{pmatrix}.
\]
中心元 \(-1\) 作用为 \(-I_U\)，故 \(U|_N\cong W\oplus W\)，而 \(E=\operatorname{Hom}_N(W,U)\) 经 \(f\mapsto f(1)\) 识别为 \(\mathbb C^2\)。继续采用 \(T_q=1\)，重数空间上的公式成为
\[
S(a,b)=\pi(i)^a\pi(j)^b.
\]
其乘积产生的符号正是上面算出的 \(\alpha\)，而 \(\alpha^{-1}=\alpha\)，符合定理中的逆因子集。矩阵 \(\pi(i)\) 的两条特征直线被 \(\pi(j)\) 交换，因此没有一条直线被全部 \(S(a,b)\) 保持，这个射影表示不可约。商群 \(\mathbb F_2^2\) 虽然交换，它的这一射影表示却是二维的；这个维数正是 \(W\) 在 \(U|_N\) 中的重数。
这里障碍余圈的值为 \(\pm1\)，故 \([\alpha]^2=1\)，而前面已由不对称性证明 \([\alpha]\ne1\)，所以其阶恰为 \(2\)。因此每个限制为 \(W^{\oplus e}\) 的非零 \(Q_8\)-表示都必须有偶数重数，当前的二维表示达到最小可能的 \(e=2\)。


\begin{chaptersummary}
双陪集 \(K\backslash G/H\) 描述 \(K\) 在 \(G/H\) 上的轨道，交集 \(K\cap xHx^{-1}\) 则是相应点的稳定子群。Mackey 分解把这一集合分解提升为任意域上的模同构；配合诱导与余诱导的伴随性质，便可把两个诱导表示之间的交织空间化为交集子群上的计算。在复数域上，自同态空间的维数又给出不可约性判据。

\ProseParagraphBreak

当 \(N\) 正规时，一个不可约 \(G\)-表示限制到 \(N\) 后，只出现一个共轭轨道，轨道中各类型的重数相等。Clifford 对应说明，其中一个等型分量是惯性群的不可约表示，诱导它就能恢复原表示。惯性群固定不可约类型，却未必允许在原空间上延拓作用；这种差别由商群上的障碍余圈记录。它膨胀到惯性群后，是原表示的射影延拓的因子集；重数空间上的射影表示以该余圈的倒数为因子集，使两个射影作用的标量相互抵消，最终形成普通表示。
\end{chaptersummary}

\BookExercises

\begin{exercise}\label{b5:ex:04-normal-double-cosets}
设 \(G\) 为有限群，\(H\trianglelefteq G\)，\(K\leq G\)。证明 \(KH\) 是子群，且每个双陪集 \(KxH\) 等于右陪集 \(KHx\)。求双陪集个数以及每块的大小，并解释为什么所有交集子群相同。
\end{exercise}
\begin{solution}
正规性给出
\((k_1h_1)(k_2h_2)=k_1k_2(k_2^{-1}h_1k_2)h_2\in KH\)，逆元也可写成 \(K\) 中元素乘 \(H\) 中元素，所以 \(KH\leq G\)。由于 \(xHx^{-1}=H\)，有 \(KxH=KHx\)。这实际上是 \(KH\) 的右陪集，因而双陪集个数为 \([G:KH]\)，每块大小为 \(|KH|=|K|\cdot|H|/|K\cap H|\)。交集 \(K\cap xHx^{-1}=K\cap H\) 不依赖 \(x\)。
\end{solution}

\begin{exercise}\label{b5:ex:04-dihedral-double-cosets}
在 \(D_8=\langle r,s:r^4=s^2=1,\ srs=r^{-1}\rangle\) 中取 \(H=\langle s\rangle\)。求 \(H\backslash D_8/H\)。若 \(\varepsilon(s)=-1\)，把 \(\operatorname{Res}_H^{D_8}\operatorname{Ind}_H^{D_8}\varepsilon\) 分解为不可约复 \(H\)-表示。
\end{exercise}
\begin{solution}
三个双陪集为 \(H\)、\(HrH=\{r,r^3,rs,r^3s\}\)、\(Hr^2H=r^2H\)，大小为 \(2,4,2\)。代表元 \(1,r,r^2\) 对应的交集分别为 \(H,\{1\},H\)，因为 \(rsr^{-1}=r^2s\)，而 \(r^2\) 在中心。Mackey 分解给出
\[
\varepsilon\oplus\operatorname{Ind}_{\{1\}}^H\mathbb C\oplus\varepsilon
\cong\mathbf1_H\oplus\varepsilon^{\oplus3}.
\]
中间一项是二阶群的正则表示，其两个特征向量分别为两个基向量之和与之差。
\end{solution}

\begin{exercise}\label{b5:ex:04-unequal-subgroups}
令 \(G=S_4\)，\(H\) 为固定 \(4\) 的子群，\(K=\bigl\langle(12)\bigr\rangle\)。求 \(K\backslash G/H\) 的代表元和交集子群，并计算
\(\dim_k\operatorname{Hom}_{k[G]}\bigl(k[G/K],k[G/H]\bigr)\)，其中 \(k\) 是任意域。
\end{exercise}
\begin{solution}
由 \(gH\mapsto g(4)\) 把 \(G/H\) 识别为四点集；\(K=\langle(12)\rangle\) 只交换 \(1,2\)，所以\(K\) 的轨道为 \(\{1,2\}\)、\(\{3\}\)、\(\{4\}\)。可取代表元 \((14),(34),1\)。第一点的稳定子群为 \(\{1\}\)，后两点的稳定子群为 \(K\)，由\cref{b5:prop:double-coset-orbits}，点 \(xH\) 的 \(K\)-稳定子群为 \(K\cap xHx^{-1}\)，所以这些也是对应的交集子群。对\cref{b5:thm:mackey-intertwining}取两个平凡表示，每个直和项均为 \(\operatorname{Hom}_{k[L_x]}(k,k)=k\)，故维数为 \(3\)，不受 \(k\) 的特征影响。
\end{solution}

\begin{exercise}\label{b5:ex:04-modular-mackey}
设 \(k\) 的特征为 \(3\)，\(G=S_3\)，\(H=\bigl\langle(12)\bigr\rangle\)，\(N=A_3\)。用 Mackey 公式计算 \(\operatorname{Res}_N^G\operatorname{Ind}_H^G k\)，并证明所得 \(N\)-模不是完全可约的。
\end{exercise}
\begin{solution}
\(N\) 在三点集 \(G/H\) 上传递，且稳定子群为 \(N\cap H=\{1\}\)。因此只有一个双陪集，限制表示同构于 \(k[N]\)。若 \(a=(123)\)，则
\[
k[N]\cong k[t]/(t^3-1)=k[t]/(t-1)^3.
\]
在正则模上，乘 \(t-1\) 把 \(1\) 送到非零的剩余类 \(t-1\)，而它的三次幂为零。在任何非零简单 \(k[N]\)-模上，这个幂零算子的核也必非零：若它单射，则其三次幂仍单射，不能是零算子。又因 \(k[N]\) 交换，这个核被群代数作用保持；简单性使核等于整个模，所以 \(t-1\) 在该模上作用为零。若正则模完全可约，它在各个简单直和项上、从而在整个正则模上都作用为零，与它把 \(1\) 送到 \(t-1\ne0\) 矛盾。这说明 Mackey 同构本身并不以半单性为前提。
\end{solution}

\begin{exercise}\label{b5:ex:04-trivial-intersection-obstruction}
设 \(G\) 为有限群，\(H<G\)。若存在 \(x\notin H\) 使 \(H\cap xHx^{-1}=\{1\}\)，证明对每个非零有限维复 \(H\)-表示 \(W\)，\(\operatorname{Ind}_H^G W\) 都可约。
\end{exercise}
\begin{solution}
在交织空间公式中，\(x=1\) 的项至少一维，所给 \(x\) 的项是底层向量空间之间的 \(\operatorname{Hom}_{\mathbb C}(W,W)\)，维数为 \((\dim W)^2>0\)。所以诱导表示的自同态空间维数至少为 \(2\)。它是完全可约的，而\cref{b5:prop:isotypic-decomposition}给出自同态空间的维数为各重数的平方和；不可约情形这个数应为 \(1\)，故该诱导表示可约。这里不必事先假设 \(W\) 不可约。
\end{solution}

\begin{exercise}\label{b5:ex:04-normal-induction-end}
设 \(G\) 为有限群，\(N\trianglelefteq G\)，\(W\) 为不可约复 \(N\)-表示，\(I=I_G(W)\)。证明
\[
\dim_{\mathbb C}\operatorname{End}_G
\left(\operatorname{Ind}_N^G W\right)=[I:N].
\]
若 \(W'\) 是另一个不可约复 \(N\)-表示，证明 \(\operatorname{Ind}_N^G W\) 与 \(\operatorname{Ind}_N^G W'\) 之间存在非零同态，当且仅当 \(W'\) 与 \(W\) 共轭。
\end{exercise}
\begin{solution}
双陪集 \(N\backslash G/N\) 就是 \(G/N\)，每个交集都是 \(N\)。交织公式中的项为 \(\operatorname{Hom}_N(W,{}^gW)\)，由 Schur 引理，它在 \(g\in I\) 时一维，在其他情形为零。故总维数为 \([I:N]\)。对两个不同的表示应用同一公式，直和项为 \(\operatorname{Hom}_N(W,{}^gW')\)，其非零条件恰为 \(W\cong{}^gW'\)，得到第二个断言。
\end{solution}

\begin{exercise}\label{b5:ex:04-nonnormal-restriction}
设 \(V=\left\{(z_1,z_2,z_3,z_4)\in\mathbb C^4:\sum_i z_i=0\right\}\)，赋予 \(S_4\) 的坐标置换作用。证明 \(V\) 不可约，而它限制到固定 \(4\) 的 \(S_3\) 后出现维数不同的不可约成分。说明这与 Clifford 限制定理没有矛盾。
\end{exercise}
\begin{solution}
\(S_4\) 在四点上的置换表示包含一次平凡表示，且对 \(H=S_3\) 使用\cref{b5:thm:mackey-intertwining}，双陪集代表元为 \(1,(34)\)，两个平凡表示的交织空间各一维，故自同态空间维数为 \(2\)。有 \(\mathbb C^4=\mathbf1\oplus V\)，而 \(V\) 中的固定向量必须各坐标相等，再由坐标和为零得到它只能是零向量。因此 \(V\) 没有平凡成分。由完全可约性，写 \(V=\bigoplus_i S_i^{\oplus m_i}\)，其中 \(S_i\) 为互不同构的非平凡不可约表示。由\cref{b5:prop:isotypic-decomposition}，
\[
2=\dim_{\mathbb C}\operatorname{End}_{S_4}(\mathbb C^4)
=1+\sum_i m_i^2.
\]
故只有一个 \(m_i=1\)，其余重数为零，\(V\) 不可约。限制到 \(S_3\) 时，有分解
\[
V=\mathbb C(1,1,1,-3)
\oplus\bigl\{(z_1,z_2,z_3,0):z_1+z_2+z_3=0\bigr\}.
\]
前一项平凡，后一项是三点置换表示去掉平凡项后的二维不可约表示，这里 \(S_2\) 在三点集上的轨道是一个固定点及另两点，故有两个双陪集。Mackey 交织公式给出三点置换表示的自同态维数为 \(2\)；去掉重数为一的平凡项后，余下各不可约重数的平方和为 \(1\)，所以这二维项不可约。这个 \(S_3\) 不是 \(S_4\) 的正规子群，所以不可约成分不必形成同维的共轭轨道。
\end{solution}

\begin{exercise}\label{b5:ex:04-s4-klein-four}
令 \(N=\{1,a,b,c\}\trianglelefteq S_4\)，其中
\(a=(12)(34)\)、\(b=(13)(24)\)、\(c=(14)(23)\)。取上题的三维表示 \(V\)，并令 \(\lambda(a)=1\)、\(\lambda(b)=\lambda(c)=-1\)。求 \(V_\lambda\) 和 \(I_{S_4}(\lambda)\)，并把 \(V\) 写成一个一维表示的诱导。
\end{exercise}
\begin{solution}
向量 \(v=(1,1,-1,-1)\) 满足 \(av=v\)、\(bv=cv=-v\)，故 \(\mathbb Cv\subseteq V_\lambda\)。对另两个分拆 \(13|24\)、\(14|23\) 取对应的正负向量，得到三条线性无关的特征直线，因此 \(V_\lambda=\mathbb Cv\)，其余两条是另外两个非平凡特征标的等型分量。

惯性群必须固定 \(\ker\lambda=\{1,a\}\)，等价于在共轭下固定 \(a\)。它就是保持无序分拆 \(\bigl\{\{1,2\},\{3,4\}\bigr\}\) 的子群 \(I\)，阶为 \(8\)。\(I\) 在 \(\mathbb Cv\) 上的特征标 \(\eta\) 对保持两块的元素取 \(1\)，对交换两块的元素取 \(-1\)。Clifford 对应给出 \(V\cong\operatorname{Ind}_I^{S_4}\eta\)，维数为 \([S_4:I]=3\)。
\end{solution}

\begin{exercise}\label{b5:ex:04-quaternion-restrictions}
在 \(\mathbb C^2\) 上令
\[
\rho(i)=\begin{pmatrix}\mathrm i&0\\0&-\mathrm i\end{pmatrix},
\qquad
\rho(j)=\begin{pmatrix}0&1\\-1&0\end{pmatrix},
\]
其中 \(\mathrm i^2=-1\)。证明它们给出 \(Q_8\) 的不可约表示，并分别对正规子群 \(\langle i\rangle\) 和 \(\{\pm1\}\) 计算 Clifford 轨道、惯性群和重数。
\end{exercise}
\begin{solution}
两个矩阵的平方都是 \(-I_2\)，且反对易，因此满足 \(Q_8\) 的关系。第一个矩阵的两个特征值不同，任何共同不变直线只能是其一条坐标特征线，但第二个矩阵交换这两条线，所以表示不可约。

对 \(\langle i\rangle\)，限制为特征标 \(i\mapsto\mathrm i\) 与 \(i\mapsto-\mathrm i\) 之和。\(j\) 交换二者，惯性群为 \(\langle i\rangle\)，重数 \(e=1\)。对中心 \(\{\pm1\}\)，只有特征标 \(-1\mapsto-1\)，出现重数 \(e=2\)，其惯性群为全部 \(Q_8\)。同一个表示的轨道长度和重数会随正规子群的选择而改变。
\end{solution}

\begin{exercise}\label{b5:ex:04-central-subgroup}
设 \(G\) 为有限群，\(N\leq Z(G)\)，\(V\) 为不可约复 \(G\)-表示。证明存在唯一一维特征标 \(\lambda:N\to\mathbb C^\times\)，使 \(n\) 在 \(V\) 上按 \(\lambda(n)I_V\) 作用。给出它与 Clifford 限制公式中 \(e\) 和 \(I\) 的关系。
\end{exercise}
\begin{solution}
每个 \(\rho(n)\) 与整个 \(\rho(G)\) 对易，故\cref{b5:cor:schur-algebraically-closed}使它为唯一标量 \(\lambda(n)\)。群乘法说明 \(\lambda(nm)=\lambda(n)\lambda(m)\)，所以它是一维特征标。限制为 \(\lambda^{\oplus\dim V}\)，从而 \(I=G\)、\(e=\dim V\)。中心上的同构类型虽然只有一个，重数仍可大于 \(1\)，上题的 \(Q_8\) 就给出例子。
\end{solution}

\begin{exercise}\label{b5:ex:04-extension-torsor}
设 \(N\trianglelefteq I\) 为有限群的正规子群，\(\rho\) 为不可约复 \(N\)-表示，并已有一个延拓 \(\widetilde\rho\)。证明所有延拓都唯一写成
\[
g\longmapsto\lambda(gN)\widetilde\rho(g),
\]
其中 \(\lambda:I/N\to\mathbb C^\times\) 为一维特征标。
\end{exercise}
\begin{solution}
给定另一延拓 \(\sigma\)，\(\sigma(g)\) 与 \(\widetilde\rho(g)\) 实现同一个 \(N\) 的共轭作用，因此它们的比值与 \(\rho(N)\) 对易。Schur 引理给出唯一标量 \(\lambda(g)\)，使 \(\sigma(g)=\lambda(g)\widetilde\rho(g)\)。比较 \(gh\) 上的乘法得到 \(\lambda(gh)=\lambda(g)\lambda(h)\)，在 \(N\) 上两延拓相同使 \(\lambda|_N=1\)，故它下降为 \(I/N\) 的特征标。反过来，用任意这样的特征标相乘仍是延拓，且表达式唯一。
\end{solution}

\begin{exercise}\label{b5:ex:04-cyclic-quotient-extension}
设 \(I\) 为有限群，\(N\trianglelefteq I\)，且 \(I/N\) 为阶 \(d\) 的循环群。证明每个同构类被 \(I\) 固定的不可约复 \(N\)-表示都能延拓到 \(I\)。
\end{exercise}
\begin{solution}
选取 \(t\in I\)，使 \(tN\) 生成 \(I/N\)，并选可逆 \(T\) 满足 \(T\rho(n)T^{-1}=\rho(tnt^{-1})\)。\(T^d\) 与 \(\rho(t^d)\) 实现同一个共轭，所以 \(T^d=a\cdot\rho(t^d)\)，其中 \(a\in\mathbb C^\times\)。取 \(b\in\mathbb C^\times\) 满足 \(b^d=a^{-1}\)，令 \(S=bT\)，便有 \(S^d=\rho(t^d)\)。

每个群元素唯一写成 \(nt^j\)，其中 \(0\leq j<d\)。定义 \(\widetilde\rho(nt^j)=\rho(n)S^j\)。乘法中先用 \(S^j\rho(m)S^{-j}=\rho(t^jmt^{-j})\) 移动 \(N\) 中的因子；若指数越过 \(d\)，再以 \(S^d=\rho(t^d)\) 作相同的归并。于是 \(\widetilde\rho(gh)=\widetilde\rho(g)\widetilde\rho(h)\)，且它在 \(N\) 上等于 \(\rho\)。复数中 \(d\) 次根的存在是这一步调整所用的条件。

这两类检查足够，是因为 \(I\) 由 \(N\) 和 \(t\) 生成，\(N\) 内部的乘法已经由 \(\rho\) 保持；将任意乘积整理成上述唯一形式，只需把 \(t\) 移过 \(N\) 中元素，并将 \(t^d\) 归并到 \(N\) 中。这里 \(t^d\in N\) 并不要求 \(t^d=1\)。例如 \(I=C_4=\langle t\rangle\)、\(N=\langle t^2\rangle\)，若 \(\rho(t^2)=-1\)，则商群虽为 \(C_2\)，延拓却必须满足 \(\widetilde\rho(t)^2=-1\)。取 \(\widetilde\rho(t)=\mathrm i\) 或 \(-\mathrm i\) 均可，而取 \(1\) 则不满足原群的乘法。
\end{solution}

\begin{exercise}\label{b5:ex:04-invariant-needed}
设 \(G=S_3\)，\(N=A_3\)，\(\lambda((123))=\zeta\)，其中 \(\zeta\neq1\) 为三次单位根。证明 \(\lambda\) 不能延拓到 \(G\)，并说明上一题的循环商群条件为何没有给出矛盾。
\end{exercise}
\begin{solution}
换位把 \((123)\) 共轭为其逆，因此把 \(\lambda\) 变成不同的 \(\lambda^{-1}\)。若 \(\lambda\) 有一维延拓，则共轭在一维标量矩阵上不起作用，迫使 \(\zeta=\zeta^{-1}\)，矛盾。虽然 \(G/N\cong C_2\) 为循环群，但上一题还要求不可约类型被整个 \(G\) 固定；这里惯性群只有 \(N\)。
\end{solution}

\begin{exercise}\label{b5:ex:04-twisted-klein-algebra}
令 \(Q=\mathbb F_3^2\)，取本原三次单位根 \(\zeta\)，并定义
\(\alpha((a,b),(c,d))=\zeta^{bc}\)。证明 \(\alpha\) 是归一化 \(2\)-余圈但不是余边。再令
\[
 X=\operatorname{diag}(1,\zeta,\zeta^2),\qquad
 Y=\begin{pmatrix}0&1&0\\0&0&1\\1&0&0\end{pmatrix}.
\]
证明 \(u_{(a,b)}\mapsto X^aY^b\) 给出 \(\mathbb C^\alpha[Q]\cong M_3(\mathbb C)\)，并求相应不可约射影表示的同构类型和维数。所有指数按模 \(3\) 计算。
\end{exercise}
\begin{solution}
双线性指数给出余圈恒等式，单位元处的值为 \(1\)。又 \(\alpha((1,0),(0,1))=1\)，而 \(\alpha((0,1),(1,0))=\zeta\ne1\)，所以由交换群的余边对称性，它不是余边。

在基 \(e_0,e_1,e_2\) 下，\(Xe_j=\zeta^j e_j\)、\(Ye_j=e_{j-1}\)，下标按模 \(3\) 计算。故 \(YX=\zeta XY\) 且 \(X^3=Y^3=I_3\)，从而
\[
 (X^aY^b)(X^cY^d)=\zeta^{bc}X^{a+c}Y^{b+d}.
\]
这证明映射保持乘法。固定 \(b\) 时，三个矩阵 \(X^aY^b\) 的非零位置相同；若 \(\sum_{a=0}^2c_aX^aY^b=0\)，作用于 \(e_j\) 后得到 \(\sum_a c_a\zeta^{a(j-b)}=0\)。当 \(j\) 遍历 \(0,1,2\) 时，\(j-b\) 也遍历这三个剩余类；所得系数矩阵 \((\zeta^{ai})_{0\le i,a\le2}\) 是以 \(1,\zeta,\zeta^2\) 为节点的 Vandermonde 矩阵，其行列式非零，故各 \(c_a=0\)。不同 \(b\) 的非零位置互不相交，故全部九个矩阵线性无关，构成 \(M_3(\mathbb C)\) 的一组基。

若 \(E\) 是非零简单 \(M_3(\mathbb C)\)-模，先由 \(\sum_iE_{ii}=1\) 取 \(0\ne w\in E_{ii}E\)，再取 \(v=E_{1i}w\ne0\)，因为 \(E_{i1}E_{1i}w=w\)。于是三个向量 \(E_{j1}v\) 均非零，分别落在不同的 \(E_{jj}E\) 中，因而线性无关。它们张成的子空间被所有矩阵单位保持，由简单性等于 \(E\)，其作用就是标准三维矩阵作用。因此相应不可约射影表示只有一个同构类型，维数为 \(3\)。
\end{solution}

\begin{exercise}\label{b5:ex:04-quaternion-factor-change}
设 \(Q=\mathbb F_2^2\) 为加法群，\(\alpha((a,b),(c,d))=(-1)^{bc}\) 为复标量因子集。比较 \(\alpha\) 与 \(Q_8\) 延拓障碍产生的余圈
\(\alpha'\bigl((a,b),(c,d)\bigr)=(-1)^{ac+bd+bc}\)，各指数按模 \(2\) 计算。求一个显式函数 \(f:Q\to\mathbb C^\times\)，使 \(\alpha'=(\delta f)\alpha\)，其中 \((\delta f)(p,q)=f(p)f(q)/f(p+q)\)。解释为何两者描述同一个非零障碍类。
\end{exercise}
\begin{solution}
取 \(f(a,b)=\mathrm i^{a+b}\)，其中 \(a,b\) 用 \(0,1\) 表示。模 \(2\) 加法发生进位时，
\[
\frac{f(a,b)f(c,d)}{f(a+c,b+d)}
=(-1)^{ac+bd}.
\]
分母的参数按模 \(2\) 相加。因此 \(\alpha'=(\delta f)\alpha\)，它们相差余边，代表同一个类。这个类非零，因为 \(\alpha((1,0),(0,1))=1\)，反向为 \(-1\)，而交换群上的余边对称。相应扭曲群代数通过重标定基同构，且\cref{b5:prop:twisted-commutator-center}说明它们均同构于 \(M_2(\mathbb C)\)。
\end{solution}

\begin{exercise}\label{b5:ex:04-factor-determinant}
设 \(Q\) 为有限群，\(\alpha\) 为归一化复标量 \(2\)-余圈。在以 \(e_q\)、\(q\in Q\) 为基的向量空间上，定义 \(R(p)e_q=\alpha(p,q)e_{pq}\)。证明 \(R\) 是因子集为 \(\alpha\) 的射影表示，并求 \(\det R(p)\)。再对函数 \(b(p)=\prod_{q\in Q}\alpha(p,q)\) 直接证明 \(\delta b=\alpha^{|Q|}\)，由此说明每个这样的余圈类都有有限阶，且其阶整除 \(|Q|\)。
\end{exercise}
\begin{solution}
余圈恒等式给出
\[
 R(p)R(q)e_r
 =\alpha(q,r)\alpha(p,qr)e_{pqr}
 =\alpha(p,q)R(pq)e_r.
\]
每个 \(R(p)\) 都是带非零系数的置换矩阵，因而可逆，且 \(R(1)=1\)。记 \(\lambda_p:Q\to Q\)、\(q\mapsto pq\)，则
\[
 \det R(p)=\operatorname{sgn}(\lambda_p)
 \prod_{q\in Q}\alpha(p,q).
\]
这些算子也就是扭曲群代数在自身上的左乘作用。维数为 \(|Q|\)，故\cref{b5:cor:projective-dimension-obstruction}给出类的阶整除 \(|Q|\)。也可不经过行列式，固定 \(p,q\)，将余圈恒等式对全部 \(r\in Q\) 相乘：
\[
 \alpha(p,q)^{|Q|}b(pq)
 =\left(\prod_r\alpha(q,r)\right)
  \left(\prod_r\alpha(p,qr)\right)
 =b(q)b(p),
\]
其中 \(r\mapsto qr\) 是双射。由于 \(b(1)=1\)，这正是 \(\delta b=\alpha^{|Q|}\) 的显式余边表达式。
\end{solution}

\begin{exercise}\label{b5:ex:04-central-extension-factor}
设有限群 \(Q\) 上的归一化余圈 \(\alpha\) 取值于 \(m\) 次单位根群 \(\mu_m\)。在 \(\mu_m\times Q\) 上定义
\[
(z,p)(w,q)=\bigl(zw\alpha(p,q),pq\bigr).
\]
证明它是一个群，且 \(\mu_m\times\{1\}\) 是中心子群。证明因子集为 \(\alpha\) 的射影表示 \(R\) 可以提升为这个群的普通表示。
\end{exercise}
\begin{solution}
三项乘法的标量部分相等，恰好等价于余圈等式；单位元为 \((1,1)\)。把余圈等式用于 \((p,p^{-1},p)\)，得到 \(\alpha(p,p^{-1})=\alpha(p^{-1},p)\)，因此逆元为
\[
(z,p)^{-1}
=\bigl(z^{-1}\cdot\alpha(p,p^{-1})^{-1},p^{-1}\bigr).
\]
归一化条件说明 \((z,1)\) 与所有元素对易；投影到 \(Q\) 的核正是这个中心子群。定义 \(\widetilde R(z,p)=z\cdot R(p)\)，则
\(\widetilde R(z,p)\widetilde R(w,q)=zw\alpha(p,q)R(pq)\)，与群乘法相容。这一表示把中心元素 \((z,1)\) 送到 \(zI\)。
\end{solution}

\begin{exercise}\label{b5:ex:04-twisted-center}
取加法群 \(Q=(\mathbb Z/4\mathbb Z)^2\)，令 \(\alpha((a,b),(c,d))=(-1)^{bc}\)。验证它是归一化 \(2\)-余圈，求 \(A=\mathbb C^\alpha[Q]\) 的中心，并给出中心的四个两两正交的非零幂等元之和为 \(1\) 的分解。再利用这些幂等元求 \(A\) 的矩阵代数分解及不可约射影表示的维数和同构类型数。
\end{exercise}
\begin{solution}
指数的奇偶性在模 \(4\) 下良定义，且双线性使余圈恒等式成立。\cref{b5:prop:twisted-commutator-center}中的对易标量仍为 \((-1)^{bc+ad}\)。取 \((c,d)=(1,0),(0,1)\)，中心条件恰好要求 \(a,b\) 均为偶数。因此
\[
 Z(A)=\operatorname{span}_{\mathbb C}
 \{u_{(0,0)},u_{(2,0)},u_{(0,2)},u_{(2,2)}\}.
\]
令 \(x=u_{(1,0)}\)、\(y=u_{(0,1)}\)，则 \(x^4=y^4=1\)、\(yx=-xy\)，中心由两个平方为 \(1\) 的元素 \(x^2,y^2\) 生成。所求幂等元为
\[
 e_{\lambda,\mu}
 =\frac14(1+\lambda x^2)(1+\mu y^2)
 \qquad(\lambda,\mu\in\{1,-1\}).
\]
两个平方为 \(1\) 的中心元素分别给出正交的二项投影，所以这四个幂等元正交且和为 \(1\)。对固定的 \(\lambda,\mu\)，四个元素 \(e_{\lambda,\mu}x^ay^b\)、\(a,b\in\{0,1\}\)，在群元素基上的支集分别具有不同奇偶性，故均非零且线性无关；高次幂由 \(x^2e_{\lambda,\mu}=\lambda e_{\lambda,\mu}\)、\(y^2e_{\lambda,\mu}=\mu e_{\lambda,\mu}\) 约化，所以它们构成 \(Ae_{\lambda,\mu}\) 的一组基。

取 \(\ell^2=\lambda\)、\(m^2=\mu\)，并取\cref{b5:prop:twisted-commutator-center}中的矩阵 \(X,Y\)。令 \(xe_{\lambda,\mu}\mapsto\ell X\)、\(ye_{\lambda,\mu}\mapsto mY\)、\(e_{\lambda,\mu}\mapsto I_2\)。两个矩阵的平方及反对易关系与这个分块相同，四个基元素的像仍线性无关，故每块同构于 \(M_2(\mathbb C)\)。于是
\[
 A\cong\prod_{\lambda,\mu\in\{1,-1\}}M_2(\mathbb C).
\]
简单模上四个中心幂等元恰有一个作用为恒等，其余为零；每块又只有一个二维简单模类型。因此相应不可约射影表示共有四个同构类型，全部为二维。
\end{solution}

\begin{exercise}\label{b5:ex:04-direct-product-clifford}
设 \(N,Q\) 为有限群，\(G=N\times Q\)。利用本章的重数空间构造，证明 \(G\) 的每个不可约复表示都同构于 \(W\otimes E\)，其中 \(W,E\) 分别是 \(N,Q\) 的不可约复表示，且两者的同构类型唯一；作用为 \((n,q)(w\otimes e)=nw\otimes qe\)。
\end{exercise}
\begin{solution}
把 \(N\) 识别为第一因子。它的每个不可约类型均被 \(G\) 固定，因为第二因子与它对易，第一因子的共轭只改变同构表示。Clifford 限制定理使任何不可约 \(G\)-表示限制为 \(W^{\oplus e}\)，且 \(W\) 的同构类型唯一。用 \((n,q)\mapsto\rho_W(n)\) 延拓 \(W\)，可在商群 \(Q\) 上取所有 \(T_q=1\)，障碍余圈为 \(1\)。\cref{b5:thm:clifford-projective-multiplicity}遂给出普通 \(Q\)-表示 \(E=\operatorname{Hom}_N(W,V)\)，并说明 \(V\) 不可约当且仅当 \(E\) 不可约。求值同构给出所述张量积，也保证 \(E\) 的同构类型唯一。反过来，每一对不可约 \(W,E\) 都由该定理给出不可约张量积。
\end{solution}

\begin{exercise}\label{b5:ex:04-clifford-square-sum}
设 \(G\) 为有限群，\(N\trianglelefteq G\)，\(W\) 为不可约复 \(N\)-表示，\(I=I_G(W)\)。让 \(U_1,\ldots,U_r\) 遍历限制到 \(N\) 时含有 \(W\) 的不可约 \(I\)-表示，写 \(U_j|_N\cong W^{\oplus e_j}\)。证明
\[
\sum_{j=1}^r e_j^2=[I:N].
\]
解释它在障碍余圈非平凡时说明了什么。
\end{exercise}
\begin{solution}
把完全可约的 \(\operatorname{Ind}_N^I W\) 分解为不可约表示。由 Frobenius 互反性，\(U_j\) 的重数为
\(\dim\operatorname{Hom}_N(W,U_j|_N)=e_j\)，其他不可约表示的重数为零。因此
\[
[I:N]\dim W
=\dim\operatorname{Ind}_N^I W
=\sum_j e_j\dim U_j
=\sum_j e_j^2\dim W.
\]
消去非零的 \(\dim W\) 得到公式。若 \(\alpha\) 为与 \(W\) 对应的障碍余圈，由重数空间定理，\(e_j\) 是因子集为 \(\alpha^{-1}\) 的全部不可约射影表示的维数。因此即使障碍非零，仍有相应的平方和公式；不能把这些射影表示直接换成普通商群表示。
\end{solution}

\begin{exercise}\label{b5:ex:04-schurian-not-simple}
设 \(k\) 的特征为 \(3\)，让 \(S_3\) 在
\[
A=\bigl\{(x,y,z)\in k^3:x+y+z=0\bigr\}
\]
上置换坐标。证明 \(A\) 可约，但 \(\operatorname{End}_{k[S_3]}A=k\)。由此说明 Mackey 不可约性判据的证明为何使用完全可约性。
\end{exercise}
\begin{solution}
非零向量 \((1,1,1)\) 属于 \(A\)，其生成的直线被全群固定，所以二维空间 \(A\) 可约。取基 \(u=(1,-1,0)\)、\(v=(0,1,-1)\)。在此基下，\((123)\) 与 \((12)\) 的矩阵分别为
\[
C=\begin{pmatrix}0&-1\\1&-1\end{pmatrix},
\qquad
T=\begin{pmatrix}-1&1\\0&1\end{pmatrix}.
\]
若算子 \(F\) 与 \(C\) 对易，写 \(F(u)=au+bv\)，则 \(F(v)=F(Cu)=CF(u)=-bu+(a-b)v\)，故 \(F=aI+bC\)。它还与 \(T\) 对易，当且仅当 \(b(CT-TC)=0\)。矩阵 \(CT-TC\) 非零，例如在 \(u\) 上的值为 \(-u-2v\neq0\)，于是 \(b=0\)。所以所有群模自同态都是标量。自同态空间一维在缺少完全可约性时，并不能排除真子表示。
\end{solution}

\begin{exercise}\label{b5:ex:04-order-twenty-one}
设 \(G=\langle r,s:r^7=s^3=1,\ srs^{-1}=r^2\rangle\)，并令 \(N=\langle r\rangle\)。求 \(N\) 的非平凡一维复特征标在 \(G\) 共轭下的轨道，构造由它们得到的不可约三维表示，并证明这些表示与三个一维表示穷尽 \(G\) 的不可约复表示。
\end{exercise}
\begin{solution}
取本原七次单位根 \(\zeta\)，写 \(\lambda_j(r)=\zeta^j\)。由于 \(s^{-1}rs=r^4\)，共轭把 \(j\) 送到 \(4j\)。六个非零指数分成 \(\{1,2,4\}\) 与 \(\{3,5,6\}\) 两个轨道，各特征标的惯性群都是 \(N\)。Clifford 对应给出两个三维不可约表示 \(\operatorname{Ind}_N^G\lambda_1\)、\(\operatorname{Ind}_N^G\lambda_3\)；它们限制到 \(N\) 的类型不同，故不同构。

若一个不可约 \(G\)-表示限制到 \(N\) 含有平凡特征标，Clifford 限制定理说明整个 \(N\) 都平凡作用，因此该表示来自 \(G/N\cong C_3\)，共有三个一维类型。若限制中含有非平凡特征标，则它属于上述两个轨道之一，且惯性群等于 \(N\)，Clifford 对应说明它只能是相应的三维诱导表示。所有情形因而已经穷尽。
\end{solution}
